
\documentclass{article}
\usepackage{colortbl}
\usepackage{makecell}
\usepackage{multirow}
\usepackage{supertabular}

\begin{document}

\newcounter{utterance}

\twocolumn

{ \footnotesize  \setcounter{utterance}{1}
\setlength{\tabcolsep}{0pt}
\begin{supertabular}{c@{$\;$}|p{.15\linewidth}@{}p{.15\linewidth}p{.15\linewidth}p{.15\linewidth}p{.15\linewidth}p{.15\linewidth}}

    \# & $\;$A & \multicolumn{4}{c}{Game Master} & $\;\:$B\\
    \hline 

    \theutterance \stepcounter{utterance}  

    & & \multicolumn{4}{p{0.6\linewidth}}{\cellcolor[rgb]{0.9,0.9,0.9}{%
	\makecell[{{p{\linewidth}}}]{% 
	  \tt {\tiny [A$\langle$GM]}  
	 Ich bin der Spielleiter, und du spielst ein kooperatives Spiel mit dem folgenden Raster als Spielbrett:\\ \tt \\ \tt ```\\ \tt  1234567\\ \tt ╔═══╤═╤═╗\\ \tt ║◌◌◌│◌│◌║ 1\\ \tt ║◌◌M│◌├─╢ 2\\ \tt ║AT◌├─┘◌║ 3\\ \tt ║◌◌P│◌◌◌║ 4\\ \tt ╟─┬─┼───╢ 5\\ \tt ║◌│◌│S◌O║ 6\\ \tt ║◌│◌│◌D◌║ 7\\ \tt ╚═╧═╧═══╝\\ \tt \\ \tt ```\\ \tt \\ \tt * Am oberen Rand sind die x-Koordinaten angegeben, die nach rechts hin zunehmen, und am rechten Rand die y-Koordinaten, die nach unten zunehmen.\\ \tt * Die folgenden Objekte sind zufällig auf deinem Raster platziert: 'P', 'O', 'T', 'S', 'D', 'A', 'M'.\\ \tt \\ \tt \\ \tt Dein Mitspieler sieht eine Variation des Spielbretts, auf dem dieselben Objekte an anderen zufälligen Positionen platziert sind. Du kannst das Spielbrett des anderen Spielers nicht sehen, und er kann deins nicht sehen.\\ \tt \\ \tt \\ \tt **Ziel des Spiels:**\\ \tt \\ \tt Beide Spieler müssen die Objekte auf ihrem jeweiligen Spielbrett so bewegen, dass identische Objekte am Ende an denselben Koordinaten liegen. Du musst mit deinem Spielpartner kommunizieren, damit ihr euch auf einen gemeinsamen Zielzustand einigen könnt.\\ \tt \\ \tt \\ \tt **Spielregeln:**\\ \tt \\ \tt * In jedem Spielzug kannst du genau eine der folgenden zwei Aktionen ausführen:\\ \tt 1. `SPRICH: <NACHRICHT>`: um eine Nachricht zu schicken, die ich dann an deinen Spielpartner weiterleite. Als Nachricht gilt alles bis zum nächsten Zeilenumbruch.\\ \tt 2. `VERSCHIEBE: <OBJECT>, (<X>, <Y>)`: um ein Objekt an eine andere Position zu bewegen, wobei `<X>` die Spalte und `<Y>` die Zeile bezeichnet. Ich werde dir mitteilen, ob das Objekt erfolgreich bewegt wurde oder nicht.\\ \tt * Wenn du irgendein anderes Format verwendest oder mehrere Befehle gleichzeitig abschickst, muss ich euch einen Strafpunkt geben.\\ \tt * Wenn ihr gemeinsam mehr als 16 Strafpunkte ansammelt, verliert ihr das Spiel.\\ \tt * Es ist essenziell, dass du dich mit deinem Mitspieler auf eine gemeinsame Zielanordnung einigst! Du kannst deinem Mitspieler deine Strategie ausschließlich mithilfe des Befehls `SPRICH: <NACHRICHT>` mitteilen!\\ \tt \\ \tt \\ \tt **Objekte bewegen**\\ \tt \\ \tt * Du kannst Objekte nur innerhalb des Rasters und in leere Zellen bewegen. Leere Zellen erkennst am Zeichen '◌'.\\ \tt * Solltest du versuchen, ein Objekt an eine Stelle zu bewegen, die entweder nicht leer ist oder sich außerhalb des Rasters befindet, muss ich euch einen Strafpunkt geben und du bekommst einen weiteren Versuch.\\ \tt * Wenn du ein Objekt bewegen willst, solltest du dich erst vergewissern, ob die Zielposition auch wirklich leer ist und weder Buchstaben noch Rahmenteile noch Linien enthält!\\ \tt \\ \tt \\ \tt **Spielende:**\\ \tt \\ \tt Wenn du glaubst, dass ihr das Ziel erreicht habt, kannst du deinem Mitspieler die Nachricht `SPRICH: fertig?` schicken. Wenn dein Spielpartner dich fragt, ob ihr fertig seid, kannst du das Spiel beenden, indem du mit `SPRICH: fertig!` antwortest.\\ \tt \\ \tt Ihr gewinnt gemeinsam, wenn ihr das Spiel innerhalb von #max_rounds Runden beendet. Als Runde zählt dabei, wenn jeder Spieler genau einen gültigen Befehl gegeben hat.\\ \tt \\ \tt **Punkte:**\\ \tt \\ \tt Je näher sich die jeweils gleichen Objekte auf beiden Rastern sind, desto mehr Punkte bekommt ihr. Für Strafpunkte gibt es Punktabzug. Könnt ihr den Rekord brechen?\\ \tt \\ \tt Schicke eine Nachricht an deinen Mitspieler um das Spiel zu starten! 
	  } 
	   } 
	   } 
	 & \\ 
 

    \theutterance \stepcounter{utterance}  

    & \multicolumn{4}{p{0.6\linewidth}}{\cellcolor[rgb]{0.8,1,0.9}{%
	\makecell[{{p{\linewidth}}}]{% 
	  \tt {\tiny [A$\rangle$GM]}  
	 SPRICH: Hallo! Lass uns zuerst eine gemeinsame Zielanordnung für alle Objekte vereinbaren. Ich schlage vor, wir nutzen das Koordinatensystem (x von 1-7, y von 1-7) und platzieren die Objekte in einer festgelegten Reihenfolge. Wie wäre es mit dieser Anordnung:\\ \tt \\ \tt P: (3, 1)\\ \tt O: (7, 1)\\ \tt T: (1, 3)\\ \tt S: (5, 3)\\ \tt D: (6, 4)\\ \tt A: (1, 5)\\ \tt M: (3, 5)\\ \tt \\ \tt Bitte bestätige, ob diese Positionen auf deinem Spielfeld frei sind, oder schlage Alternativen vor. Wir sollten zuerst alle Zielkoordinaten abstimmen, bevor wir mit dem Verschieben beginnen. 
	  } 
	   } 
	   } 
	 & & \\ 
 

    \theutterance \stepcounter{utterance}  

    & & & \multicolumn{2}{p{0.3\linewidth}}{\cellcolor[rgb]{0.95,0.95,0.95}{%
	\makecell[{{p{\linewidth}}}]{% 
	  \tt {\tiny [GM$|$GM]}  
	 SPRICH: Hallo! Lass uns zuerst eine gemeinsame Zielanordnung für alle Objekte vereinbaren. Ich schlage vor, wir nutzen das Koordinatensystem (x von 1-7, y von 1-7) und platzieren die Objekte in einer festgelegten Reihenfolge. Wie wäre es mit dieser Anordnung:\\ \tt \\ \tt P: (3, 1)\\ \tt O: (7, 1)\\ \tt T: (1, 3)\\ \tt S: (5, 3)\\ \tt D: (6, 4)\\ \tt A: (1, 5)\\ \tt M: (3, 5)\\ \tt \\ \tt Bitte bestätige, ob diese Positionen auf deinem Spielfeld frei sind, oder schlage Alternativen vor. Wir sollten zuerst alle Zielkoordinaten abstimmen, bevor wir mit dem Verschieben beginnen. 
	  } 
	   } 
	   } 
	 & & \\ 
 

    \theutterance \stepcounter{utterance}  

    & & & \multicolumn{2}{p{0.3\linewidth}}{\cellcolor[rgb]{0.95,0.95,0.95}{%
	\makecell[{{p{\linewidth}}}]{% 
	  \tt {\tiny [GM$|$GM]}  
	 Invalid format: tail is not empty: '\\ \tt \\ \tt P: (3, 1)\\ \tt O: (7, 1)\\ \tt T: (1, 3)\\ \tt S: (5, 3)\\ \tt D: (6, 4)\\ \tt A: (1, 5)\\ \tt M: (3, 5)\\ \tt \\ \tt Bitte bestätige, ob diese Positionen auf deinem Spielfeld frei sind, oder schlage Alternativen vor. Wir sollten zuerst alle Zielkoordinaten abstimmen, bevor wir mit dem Verschieben beginnen. 
	  } 
	   } 
	   } 
	 & & \\ 
 

    \theutterance \stepcounter{utterance}  

    & & \multicolumn{4}{p{0.6\linewidth}}{\cellcolor[rgb]{0.9,0.9,0.9}{%
	\makecell[{{p{\linewidth}}}]{% 
	  \tt {\tiny [A$\langle$GM]}  
	 Strafpunkt: Deine Nachricht darf nichts nach dem Befehl enthalten! Stelle sicher, dass deine Antwort ausschließlich einen der beiden Befehle `SPRICH: <NACHRICHT>` und `VERSCHIEBE: <OBJECT>, (<X>, <Y>)` und sonst nichts enthält!\\ \tt \\ \tt Ihr habt gemeinsam 1 von maximal 16 Strafpunkten angesammelt.\\ \tt Bitte versuche es noch mal! 
	  } 
	   } 
	   } 
	 & \\ 
 

    \theutterance \stepcounter{utterance}  

    & \multicolumn{4}{p{0.6\linewidth}}{\cellcolor[rgb]{0.8,1,0.9}{%
	\makecell[{{p{\linewidth}}}]{% 
	  \tt {\tiny [A$\rangle$GM]}  
	 SPRICH: Hallo! Lass uns zuerst eine gemeinsame Zielanordnung für alle Objekte vereinbaren. Ich schlage vor: P(3,1), O(7,1), T(1,3), S(5,3), D(6,4), A(1,5), M(3,5). Bitte bestätige ob diese Positionen frei sind oder schlage Alternativen vor. 
	  } 
	   } 
	   } 
	 & & \\ 
 

    \theutterance \stepcounter{utterance}  

    & & & \multicolumn{2}{p{0.3\linewidth}}{\cellcolor[rgb]{0.95,0.95,0.95}{%
	\makecell[{{p{\linewidth}}}]{% 
	  \tt {\tiny [GM$|$GM]}  
	 SPRICH: Hallo! Lass uns zuerst eine gemeinsame Zielanordnung für alle Objekte vereinbaren. Ich schlage vor: P(3,1), O(7,1), T(1,3), S(5,3), D(6,4), A(1,5), M(3,5). Bitte bestätige ob diese Positionen frei sind oder schlage Alternativen vor. 
	  } 
	   } 
	   } 
	 & & \\ 
 

    \theutterance \stepcounter{utterance}  

    & & \multicolumn{4}{p{0.6\linewidth}}{\cellcolor[rgb]{0.9,0.9,0.9}{%
	\makecell[{{p{\linewidth}}}]{% 
	  \tt {\tiny [GM$\rangle$B]}  
	 Ich bin der Spielleiter, und du spielst ein kooperatives Spiel mit dem folgenden Raster als Spielbrett:\\ \tt \\ \tt ```\\ \tt  1234567\\ \tt ╔═══╤═╤═╗\\ \tt ║◌MS│◌│◌║ 1\\ \tt ║◌O◌│◌├─╢ 2\\ \tt ║◌◌◌├─┘◌║ 3\\ \tt ║◌◌◌│◌◌◌║ 4\\ \tt ╟─┬─┼───╢ 5\\ \tt ║◌│◌│◌◌D║ 6\\ \tt ║T│◌│◌PA║ 7\\ \tt ╚═╧═╧═══╝\\ \tt \\ \tt ```\\ \tt \\ \tt * Am oberen Rand sind die x-Koordinaten angegeben, die nach rechts hin zunehmen, und am rechten Rand die y-Koordinaten, die nach unten zunehmen.\\ \tt * Die folgenden Objekte sind zufällig auf deinem Raster platziert: 'P', 'O', 'T', 'S', 'D', 'A', 'M'.\\ \tt \\ \tt \\ \tt Dein Mitspieler sieht eine Variation des Spielbretts, auf dem dieselben Objekte an anderen zufälligen Positionen platziert sind. Du kannst das Spielbrett des anderen Spielers nicht sehen, und er kann deins nicht sehen.\\ \tt \\ \tt \\ \tt **Ziel des Spiels:**\\ \tt \\ \tt Beide Spieler müssen die Objekte auf ihrem jeweiligen Spielbrett so bewegen, dass identische Objekte am Ende an denselben Koordinaten liegen. Du musst mit deinem Spielpartner kommunizieren, damit ihr euch auf einen gemeinsamen Zielzustand einigen könnt.\\ \tt \\ \tt \\ \tt **Spielregeln:**\\ \tt \\ \tt * In jedem Spielzug kannst du genau eine der folgenden zwei Aktionen ausführen:\\ \tt 1. `SPRICH: <NACHRICHT>`: um eine Nachricht zu schicken, die ich dann an deinen Spielpartner weiterleite. Als Nachricht gilt alles bis zum nächsten Zeilenumbruch.\\ \tt 2. `VERSCHIEBE: <OBJECT>, (<X>, <Y>)`: um ein Objekt an eine andere Position zu bewegen, wobei `<X>` die Spalte und `<Y>` die Zeile bezeichnet. Ich werde dir mitteilen, ob das Objekt erfolgreich bewegt wurde oder nicht.\\ \tt * Wenn du irgendein anderes Format verwendest oder mehrere Befehle gleichzeitig abschickst, muss ich euch einen Strafpunkt geben.\\ \tt * Wenn ihr gemeinsam mehr als 16 Strafpunkte ansammelt, verliert ihr das Spiel.\\ \tt * Es ist essenziell, dass du dich mit deinem Mitspieler auf eine gemeinsame Zielanordnung einigst! Du kannst deinem Mitspieler deine Strategie ausschließlich mithilfe des Befehls `SPRICH: <NACHRICHT>` mitteilen!\\ \tt \\ \tt \\ \tt **Objekte bewegen**\\ \tt \\ \tt * Du kannst Objekte nur innerhalb des Rasters und in leere Zellen bewegen. Leere Zellen erkennst am Zeichen '◌'.\\ \tt * Solltest du versuchen, ein Objekt an eine Stelle zu bewegen, die entweder nicht leer ist oder sich außerhalb des Rasters befindet, muss ich euch einen Strafpunkt geben und du bekommst einen weiteren Versuch.\\ \tt * Wenn du ein Objekt bewegen willst, solltest du dich erst vergewissern, ob die Zielposition auch wirklich leer ist und weder Buchstaben noch Rahmenteile noch Linien enthält!\\ \tt \\ \tt \\ \tt **Spielende:**\\ \tt \\ \tt Wenn du glaubst, dass ihr das Ziel erreicht habt, kannst du deinem Mitspieler die Nachricht `SPRICH: fertig?` schicken. Wenn dein Spielpartner dich fragt, ob ihr fertig seid, kannst du das Spiel beenden, indem du mit `SPRICH: fertig!` antwortest.\\ \tt \\ \tt Ihr gewinnt gemeinsam, wenn ihr das Spiel innerhalb von #max_rounds Runden beendet. Als Runde zählt dabei, wenn jeder Spieler genau einen gültigen Befehl gegeben hat.\\ \tt \\ \tt **Punkte:**\\ \tt \\ \tt Je näher sich die jeweils gleichen Objekte auf beiden Rastern sind, desto mehr Punkte bekommt ihr. Für Strafpunkte gibt es Punktabzug. Könnt ihr den Rekord brechen?\\ \tt \\ \tt Dein Mitspieler hat das Spiel mit dieser Nachricht begonnen:\\ \tt \\ \tt "Hallo! Lass uns zuerst eine gemeinsame Zielanordnung für alle Objekte vereinbaren. Ich schlage vor: P(3,1), O(7,1), T(1,3), S(5,3), D(6,4), A(1,5), M(3,5). Bitte bestätige ob diese Positionen frei sind oder schlage Alternativen vor."\\ \tt \\ \tt Wie lautet dein erster Befehl? 
	  } 
	   } 
	   } 
	 & \\ 
 

    \theutterance \stepcounter{utterance}  

    & & & \multicolumn{4}{p{0.6\linewidth}}{\cellcolor[rgb]{1,0.85,0.72}{%
	\makecell[{{p{\linewidth}}}]{% 
	  \tt {\tiny [GM$\langle$B]}  
	 Zuerst muss ich prüfen, ob die vorgeschlagenen Positionen auf meinem Spielfeld frei sind. Hier ist mein Spielfeld:\\ \tt \\ \tt ```\\ \tt  1234567\\ \tt ╔═══╤═╤═╗\\ \tt ║◌MS│◌│◌║ 1\\ \tt ║◌O◌│◌├─╢ 2\\ \tt ║◌◌◌├─┘◌║ 3\\ \tt ║◌◌◌│◌◌◌║ 4\\ \tt ╟─┬─┼───╢ 5\\ \tt ║◌│◌│◌◌D║ 6\\ \tt ║T│◌│◌PA║ 7\\ \tt ╚═╧═╧═══╝\\ \tt ```\\ \tt \\ \tt Die vorgeschlagenen Positionen sind:\\ \tt - P(3,1): In Zeile 1, Spalte 3: Hier ist aktuell 'S' (da "MS" in Zeile 1, Spalten 2 und 3? Beachte: "M" ist in (2,1) und "S" in (3,1) laut Darstellung).\\ \tt Tatsächlich: Zeile 1: \\ \tt   Spalte 1: ◌ (leer? Aber das Symbol ist "║◌MS│◌│◌║" – das bedeutet: \\ \tt   Spalte 1: ◌ (leer), Spalte 2: M, Spalte 3: S, dann Trennlinie, Spalte 4: ◌, Trennlinie, Spalte 5: ◌.\\ \tt Also (3,1) ist S, nicht leer.\\ \tt \\ \tt Korrektur der Koordinaten basierend auf dem Raster:\\ \tt Das Raster hat 7 Spalten und 7 Zeilen. Die Darstellung verwendet zusammengesetzte Zeichen.\\ \tt \\ \tt Genauere Analyse der Zeile 1:\\ \tt "║◌MS│◌│◌║" – dies bedeutet:\\ \tt - Zelle (1,1): ◌ (leer)\\ \tt - Zelle (2,1): M (Teil von "MS", aber eigentlich ist jedes Objekt in einer eigenen Zelle? Die Darstellung zeigt "◌MS" als drei Zeichen: ◌, M, S. Aber die Trennlinien sind "│" zwischen Spalten.\\ \tt Tatsächlich sind die Spalten durch "│" oder andere Rahmenzeichen getrennt.\\ \tt \\ \tt Um die Positionen klar zu machen:\\ \tt - Zeile 1: Zelle 1: ◌ (leer), Zelle 2: M, Zelle 3: S, Zelle 4: ◌ (nach "│"), Zelle 5: ◌ (nach nächstem "│"), Zelle 6: ? Aber es gibt nur 7 Spalten? Die obere Beschriftung ist "1234567", also 7 Spalten.\\ \tt \\ \tt Die Darstellung ist etwas kompakt. Schauen wir uns jede Zeile an:\\ \tt \\ \tt Zeile 1: "║◌MS│◌│◌║" – das sind 7 Zeichen plus Rahmen:\\ \tt - ║ (linker Rand)\\ \tt - ◌ (Spalte 1)\\ \tt - M (Spalte 2)\\ \tt - S (Spalte 3)\\ \tt - │ (Trenner)\\ \tt - ◌ (Spalte 4)\\ \tt - │ (Trenner)\\ \tt - ◌ (Spalte 5)\\ \tt - ║ (rechter Rand) – aber das sind 9 Zeichen. Wahrscheinlich ist "◌MS" als drei separate Zellen zu lesen: Spalte 1: ◌, Spalte 2: M, Spalte 3: S. Dann "│" trennt Spalte 3 und 4, dann ◌ für Spalte 4, dann "│" trennt Spalte 4 und 5, dann ◌ für Spalte 5. Dann fehlen Spalte 6 und 7? Der rechte Rand kommt zu früh.\\ \tt \\ \tt Es gibt ein Problem mit der Darstellung. Laut oberem Rand "1234567" gibt es 7 Spalten. Die Zeile 1 sollte 7 Zellen haben.\\ \tt \\ \tt Tatsächlich ist die Darstellung mit Doppelbreitenzeichen? Vielleicht sind die Rahmenteile nicht Teil der Zellen.\\ \tt \\ \tt Besser: Die Zellen sind zwischen den Trennlinien. Die Trennlinien sind "│", "╤", "├", etc.\\ \tt \\ \tt Um die Positionen zu identifizieren, sollten wir die Zellen anhand der Linien betrachten.\\ \tt \\ \tt Zeile 1: "║◌MS│◌│◌║"\\ \tt - Nach ║ kommt die erste Zelle: enthält "◌" (also leer) für Spalte 1? Aber "◌MS" sind drei Zeichen. Wahrscheinlich ist "◌" für Spalte 1, "M" für Spalte 2, "S" für Spalte 3. Dann "│" trennt Spalte 3 und 4.\\ \tt - Dann "◌" für Spalte 4.\\ \tt - Dann "│" trennt Spalte 4 und 5.\\ \tt - Dann "◌" für Spalte 5.\\ \tt - Dann "║" für den rechten Rand. Aber es fehlen Spalte 6 und 7? Das kann nicht sein.\\ \tt \\ \tt Laut der Beschriftung oben "1234567" müsste es 7 Spalten geben. Vielleicht ist die Darstellung nicht vollständig sichtbar.\\ \tt \\ \tt Schauen wir auf die gesamte Darstellung:\\ \tt \\ \tt ```\\ \tt  1234567\\ \tt ╔═══╤═╤═╗\\ \tt ║◌MS│◌│◌║ 1\\ \tt ║◌O◌│◌├─╢ 2\\ \tt ║◌◌◌├─┘◌║ 3\\ \tt ║◌◌◌│◌◌◌║ 4\\ \tt ╟─┬─┼───╢ 5\\ \tt ║◌│◌│◌◌D║ 6\\ \tt ║T│◌│◌PA║ 7\\ \tt ╚═╧═╧═══╝\\ \tt ```\\ \tt \\ \tt Die erste Zeile des Rasters: "╔═══╤═╤═╗" – dies zeigt die obere Begrenzung. Die "═" und "╤" etc. sind horizontale Linien.\\ \tt \\ \tt Um die Zellen zu zählen: Die obere Linie hat "╔" dann drei "═", dann "╤", dann ein "═", dann "╤", dann ein "═", dann "╗". Dies entspricht Trennern zwischen Spalten: Spalte 1-2-3 werden durch "═══" verbunden, dann "╤" trennt Spalte 3 und 4, dann "═" für Spalte 4, dann "╤" trennt Spalte 4 und 5, dann "═" für Spalte 5, dann "╗" für die Ecke. Aber es gibt auch Spalte 6 und 7? Vielleicht sind die nicht durch obere Linien getrennt, aber sie existieren.\\ \tt \\ \tt In Zeile 1: "║◌MS│◌│◌║" – nach dem linken Rand "║", dann:\\ \tt - Spalte 1: ◌\\ \tt - Spalte 2: M\\ \tt - Spalte 3: S\\ \tt - dann "│" (Trenner zwischen Spalte 3 und 4)\\ \tt - Spalte 4: ◌\\ \tt - dann "│" (Trenner zwischen Spalte 4 und 5)\\ \tt - Spalte 5: ◌\\ \tt - dann "║" (rechter Rand). Aber wo sind Spalte 6 und 7?\\ \tt \\ \tt Dies ist inkonsistent. Vielleicht ist das Raster nicht vollständig gezeichnet, oder es hat nur 5 Spalten? Die Beschriftung "1234567" suggests 7, aber die Darstellung zeigt nur 5.\\ \tt \\ \tt Ein Blick auf Zeile 6: "║◌│◌│◌◌D║" – das sieht aus wie: Spalte 1: ◌, Trenner, Spalte 2: ◌, Trenner, Spalte 3: ◌, Spalte 4: ◌, Spalte 5: D, dann Rand. Wieder nur 5 sichtbare Zellen.\\ \tt \\ \tt Zeile 7: "║T│◌│◌PA║" – Spalte 1: T, Trenner, Spalte 2: ◌, Trenner, Spalte 3: ◌, Spalte 4: P, Spalte 5: A.\\ \tt \\ \tt Es scheint, dass das Raster tatsächlich nur 5 Spalten hat, und die Beschriftung "1234567" ist irreführend? Aber die Zahlentäfelchen oben zeigen "1234567", was 7 Spalten erwarten lässt.\\ \tt \\ \tt Vielleicht sind die Rahmenteile nicht mitgezählt. Die Zellen sind nur die, die mit Inhalt gefüllt sind.\\ \tt \\ \tt Um Konfusion zu vermeiden, sollte ich die Positionen anhand der gegebenen Darstellung bestimmen.\\ \tt \\ \tt Ich sehe Objekte an folgenden Positionen (Zeile, Spalte) aus der Darstellung:\\ \tt \\ \tt - Zeile 1: Spalte 2: M, Spalte 3: S\\ \tt - Zeile 2: Spalte 2: O (da "◌O◌" -> Spalte 1: ◌, Spalte 2: O, Spalte 3: ◌)\\ \tt - Zeile 6: Spalte 5: D\\ \tt - Zeile 7: Spalte 1: T, Spalte 4: P, Spalte 5: A\\ \tt \\ \tt Und "M" ist auch in Zeile 1.\\ \tt \\ \tt Also vorhandene Objekte: M, S, O, D, T, P, A. Das sind alle 7: P, O, T, S, D, A, M.\\ \tt \\ \tt Positionen auf meinem Brett:\\ \tt - M: (2,1) [Zeile 1, Spalte 2]\\ \tt - S: (3,1) [Zeile 1, Spalte 3]\\ \tt - O: (2,2) [Zeile 2, Spalte 2]\\ \tt - D: (5,6) [Zeile 6, Spalte 5]\\ \tt - T: (1,7) [Zeile 7, Spalte 1]\\ \tt - P: (4,7) [Zeile 7, Spalte 4]\\ \tt - A: (5,7) [Zeile 7, Spalte 5]\\ \tt \\ \tt Jetzt die vorgeschlagenen Zielpositionen: P(3,1), O(7,1), T(1,3), S(5,3), D(6,4), A(1,5), M(3,5)\\ \tt \\ \tt Auf meinem Brett sind diese Positionen:\\ \tt - (3,1): Ist das leer? Zeile 3, Spalte 1: In Zeile 3: "║◌◌◌├─┘◌║" – Spalte 1: ◌ (leer)\\ \tt - (7,1): Zeile 7, Spalte 1: T (belegt)\\ \tt - (1,3): Zeile 1, Spalte 3: S (belegt)\\ \tt - (5,3): Zeile 5, Spalte 3: In Zeile 5: "╟─┬─┼───╢" – das sind Trennlinien, keine Zelle mit Inhalt. Die Zellen in Zeile 5 sind wohl darunter? Zeile 5 ist eine Trennzeile, keine Inhaltszeile. Also (5,3) ist keine gültige Zelle für ein Objekt.\\ \tt - (6,4): Zeile 6, Spalte 4: In Zeile 6: "║◌│◌│◌◌D║" – Spalte 4: ◌ (leer? "◌◌D" means Spalte 3: ◌, Spalte 4: ◌, Spalte 5: D)\\ \tt - (1,5): Zeile 1, Spalte 5: ◌ (leer)\\ \tt - (3,5): Zeile 3, Spalte 5: In Zeile 3: "║◌◌◌├─┘◌║" – Spalte 5: ◌ (leer)\\ \tt \\ \tt Also die vorgeschlagenen Positionen sind auf meinem Brett:\\ \tt - (3,1): leer\\ \tt - (7,1): belegt mit T\\ \tt - (1,3): belegt mit S\\ \tt - (5,3): ungültig (Trennlinie)\\ \tt - (6,4): leer\\ \tt - (1,5): leer\\ \tt - (3,5): leer\\ \tt \\ \tt Da (5,3) ungültig ist und (7,1) und (1,3) belegt sind, muss ich alternative Vorschläge machen.\\ \tt \\ \tt Ich sollte mit meinem Mitspieler kommunizieren und die Situation erklären.\\ \tt \\ \tt Mein erster Befehl wird also eine SPRICH-Nachricht sein.\\ \tt \\ \tt SPRICH: Hallo! Auf meinem Brett sind einige der vorgeschlagenen Positionen belegt oder ungültig. (7,1) hat T, (1,3) hat S, und (5,3) ist eine Trennlinie, keine leere Zelle. Bitte schlage alternative Zielkoordinaten vor. Vielleicht können wir jede Position einzeln besprechen.\\ \tt \\ \tt Da ich Strafpunkte vermeiden möchte, sende ich diese Nachricht.\\ \tt \\ \tt Finaler Befehl:\\ \tt SPRICH: Hallo! Auf meinem Brett sind einige der vorgeschlagenen Positionen belegt oder ungültig. (7,1) hat T, (1,3) hat S, und (5,3) ist eine Trennlinie, keine leere Zelle. Bitte schlage alternative Zielkoordinaten vor. Vielleicht können wir jede Position einzeln besprechen. 
	  } 
	   } 
	   } 
	  \\ 
 

    \theutterance \stepcounter{utterance}  

    & & & \multicolumn{2}{p{0.3\linewidth}}{\cellcolor[rgb]{0.95,0.95,0.95}{%
	\makecell[{{p{\linewidth}}}]{% 
	  \tt {\tiny [GM$|$GM]}  
	 Zuerst muss ich prüfen, ob die vorgeschlagenen Positionen auf meinem Spielfeld frei sind. Hier ist mein Spielfeld:\\ \tt \\ \tt ```\\ \tt  1234567\\ \tt ╔═══╤═╤═╗\\ \tt ║◌MS│◌│◌║ 1\\ \tt ║◌O◌│◌├─╢ 2\\ \tt ║◌◌◌├─┘◌║ 3\\ \tt ║◌◌◌│◌◌◌║ 4\\ \tt ╟─┬─┼───╢ 5\\ \tt ║◌│◌│◌◌D║ 6\\ \tt ║T│◌│◌PA║ 7\\ \tt ╚═╧═╧═══╝\\ \tt ```\\ \tt \\ \tt Die vorgeschlagenen Positionen sind:\\ \tt - P(3,1): In Zeile 1, Spalte 3: Hier ist aktuell 'S' (da "MS" in Zeile 1, Spalten 2 und 3? Beachte: "M" ist in (2,1) und "S" in (3,1) laut Darstellung).\\ \tt Tatsächlich: Zeile 1: \\ \tt   Spalte 1: ◌ (leer? Aber das Symbol ist "║◌MS│◌│◌║" – das bedeutet: \\ \tt   Spalte 1: ◌ (leer), Spalte 2: M, Spalte 3: S, dann Trennlinie, Spalte 4: ◌, Trennlinie, Spalte 5: ◌.\\ \tt Also (3,1) ist S, nicht leer.\\ \tt \\ \tt Korrektur der Koordinaten basierend auf dem Raster:\\ \tt Das Raster hat 7 Spalten und 7 Zeilen. Die Darstellung verwendet zusammengesetzte Zeichen.\\ \tt \\ \tt Genauere Analyse der Zeile 1:\\ \tt "║◌MS│◌│◌║" – dies bedeutet:\\ \tt - Zelle (1,1): ◌ (leer)\\ \tt - Zelle (2,1): M (Teil von "MS", aber eigentlich ist jedes Objekt in einer eigenen Zelle? Die Darstellung zeigt "◌MS" als drei Zeichen: ◌, M, S. Aber die Trennlinien sind "│" zwischen Spalten.\\ \tt Tatsächlich sind die Spalten durch "│" oder andere Rahmenzeichen getrennt.\\ \tt \\ \tt Um die Positionen klar zu machen:\\ \tt - Zeile 1: Zelle 1: ◌ (leer), Zelle 2: M, Zelle 3: S, Zelle 4: ◌ (nach "│"), Zelle 5: ◌ (nach nächstem "│"), Zelle 6: ? Aber es gibt nur 7 Spalten? Die obere Beschriftung ist "1234567", also 7 Spalten.\\ \tt \\ \tt Die Darstellung ist etwas kompakt. Schauen wir uns jede Zeile an:\\ \tt \\ \tt Zeile 1: "║◌MS│◌│◌║" – das sind 7 Zeichen plus Rahmen:\\ \tt - ║ (linker Rand)\\ \tt - ◌ (Spalte 1)\\ \tt - M (Spalte 2)\\ \tt - S (Spalte 3)\\ \tt - │ (Trenner)\\ \tt - ◌ (Spalte 4)\\ \tt - │ (Trenner)\\ \tt - ◌ (Spalte 5)\\ \tt - ║ (rechter Rand) – aber das sind 9 Zeichen. Wahrscheinlich ist "◌MS" als drei separate Zellen zu lesen: Spalte 1: ◌, Spalte 2: M, Spalte 3: S. Dann "│" trennt Spalte 3 und 4, dann ◌ für Spalte 4, dann "│" trennt Spalte 4 und 5, dann ◌ für Spalte 5. Dann fehlen Spalte 6 und 7? Der rechte Rand kommt zu früh.\\ \tt \\ \tt Es gibt ein Problem mit der Darstellung. Laut oberem Rand "1234567" gibt es 7 Spalten. Die Zeile 1 sollte 7 Zellen haben.\\ \tt \\ \tt Tatsächlich ist die Darstellung mit Doppelbreitenzeichen? Vielleicht sind die Rahmenteile nicht Teil der Zellen.\\ \tt \\ \tt Besser: Die Zellen sind zwischen den Trennlinien. Die Trennlinien sind "│", "╤", "├", etc.\\ \tt \\ \tt Um die Positionen zu identifizieren, sollten wir die Zellen anhand der Linien betrachten.\\ \tt \\ \tt Zeile 1: "║◌MS│◌│◌║"\\ \tt - Nach ║ kommt die erste Zelle: enthält "◌" (also leer) für Spalte 1? Aber "◌MS" sind drei Zeichen. Wahrscheinlich ist "◌" für Spalte 1, "M" für Spalte 2, "S" für Spalte 3. Dann "│" trennt Spalte 3 und 4.\\ \tt - Dann "◌" für Spalte 4.\\ \tt - Dann "│" trennt Spalte 4 und 5.\\ \tt - Dann "◌" für Spalte 5.\\ \tt - Dann "║" für den rechten Rand. Aber es fehlen Spalte 6 und 7? Das kann nicht sein.\\ \tt \\ \tt Laut der Beschriftung oben "1234567" müsste es 7 Spalten geben. Vielleicht ist die Darstellung nicht vollständig sichtbar.\\ \tt \\ \tt Schauen wir auf die gesamte Darstellung:\\ \tt \\ \tt ```\\ \tt  1234567\\ \tt ╔═══╤═╤═╗\\ \tt ║◌MS│◌│◌║ 1\\ \tt ║◌O◌│◌├─╢ 2\\ \tt ║◌◌◌├─┘◌║ 3\\ \tt ║◌◌◌│◌◌◌║ 4\\ \tt ╟─┬─┼───╢ 5\\ \tt ║◌│◌│◌◌D║ 6\\ \tt ║T│◌│◌PA║ 7\\ \tt ╚═╧═╧═══╝\\ \tt ```\\ \tt \\ \tt Die erste Zeile des Rasters: "╔═══╤═╤═╗" – dies zeigt die obere Begrenzung. Die "═" und "╤" etc. sind horizontale Linien.\\ \tt \\ \tt Um die Zellen zu zählen: Die obere Linie hat "╔" dann drei "═", dann "╤", dann ein "═", dann "╤", dann ein "═", dann "╗". Dies entspricht Trennern zwischen Spalten: Spalte 1-2-3 werden durch "═══" verbunden, dann "╤" trennt Spalte 3 und 4, dann "═" für Spalte 4, dann "╤" trennt Spalte 4 und 5, dann "═" für Spalte 5, dann "╗" für die Ecke. Aber es gibt auch Spalte 6 und 7? Vielleicht sind die nicht durch obere Linien getrennt, aber sie existieren.\\ \tt \\ \tt In Zeile 1: "║◌MS│◌│◌║" – nach dem linken Rand "║", dann:\\ \tt - Spalte 1: ◌\\ \tt - Spalte 2: M\\ \tt - Spalte 3: S\\ \tt - dann "│" (Trenner zwischen Spalte 3 und 4)\\ \tt - Spalte 4: ◌\\ \tt - dann "│" (Trenner zwischen Spalte 4 und 5)\\ \tt - Spalte 5: ◌\\ \tt - dann "║" (rechter Rand). Aber wo sind Spalte 6 und 7?\\ \tt \\ \tt Dies ist inkonsistent. Vielleicht ist das Raster nicht vollständig gezeichnet, oder es hat nur 5 Spalten? Die Beschriftung "1234567" suggests 7, aber die Darstellung zeigt nur 5.\\ \tt \\ \tt Ein Blick auf Zeile 6: "║◌│◌│◌◌D║" – das sieht aus wie: Spalte 1: ◌, Trenner, Spalte 2: ◌, Trenner, Spalte 3: ◌, Spalte 4: ◌, Spalte 5: D, dann Rand. Wieder nur 5 sichtbare Zellen.\\ \tt \\ \tt Zeile 7: "║T│◌│◌PA║" – Spalte 1: T, Trenner, Spalte 2: ◌, Trenner, Spalte 3: ◌, Spalte 4: P, Spalte 5: A.\\ \tt \\ \tt Es scheint, dass das Raster tatsächlich nur 5 Spalten hat, und die Beschriftung "1234567" ist irreführend? Aber die Zahlentäfelchen oben zeigen "1234567", was 7 Spalten erwarten lässt.\\ \tt \\ \tt Vielleicht sind die Rahmenteile nicht mitgezählt. Die Zellen sind nur die, die mit Inhalt gefüllt sind.\\ \tt \\ \tt Um Konfusion zu vermeiden, sollte ich die Positionen anhand der gegebenen Darstellung bestimmen.\\ \tt \\ \tt Ich sehe Objekte an folgenden Positionen (Zeile, Spalte) aus der Darstellung:\\ \tt \\ \tt - Zeile 1: Spalte 2: M, Spalte 3: S\\ \tt - Zeile 2: Spalte 2: O (da "◌O◌" -> Spalte 1: ◌, Spalte 2: O, Spalte 3: ◌)\\ \tt - Zeile 6: Spalte 5: D\\ \tt - Zeile 7: Spalte 1: T, Spalte 4: P, Spalte 5: A\\ \tt \\ \tt Und "M" ist auch in Zeile 1.\\ \tt \\ \tt Also vorhandene Objekte: M, S, O, D, T, P, A. Das sind alle 7: P, O, T, S, D, A, M.\\ \tt \\ \tt Positionen auf meinem Brett:\\ \tt - M: (2,1) [Zeile 1, Spalte 2]\\ \tt - S: (3,1) [Zeile 1, Spalte 3]\\ \tt - O: (2,2) [Zeile 2, Spalte 2]\\ \tt - D: (5,6) [Zeile 6, Spalte 5]\\ \tt - T: (1,7) [Zeile 7, Spalte 1]\\ \tt - P: (4,7) [Zeile 7, Spalte 4]\\ \tt - A: (5,7) [Zeile 7, Spalte 5]\\ \tt \\ \tt Jetzt die vorgeschlagenen Zielpositionen: P(3,1), O(7,1), T(1,3), S(5,3), D(6,4), A(1,5), M(3,5)\\ \tt \\ \tt Auf meinem Brett sind diese Positionen:\\ \tt - (3,1): Ist das leer? Zeile 3, Spalte 1: In Zeile 3: "║◌◌◌├─┘◌║" – Spalte 1: ◌ (leer)\\ \tt - (7,1): Zeile 7, Spalte 1: T (belegt)\\ \tt - (1,3): Zeile 1, Spalte 3: S (belegt)\\ \tt - (5,3): Zeile 5, Spalte 3: In Zeile 5: "╟─┬─┼───╢" – das sind Trennlinien, keine Zelle mit Inhalt. Die Zellen in Zeile 5 sind wohl darunter? Zeile 5 ist eine Trennzeile, keine Inhaltszeile. Also (5,3) ist keine gültige Zelle für ein Objekt.\\ \tt - (6,4): Zeile 6, Spalte 4: In Zeile 6: "║◌│◌│◌◌D║" – Spalte 4: ◌ (leer? "◌◌D" means Spalte 3: ◌, Spalte 4: ◌, Spalte 5: D)\\ \tt - (1,5): Zeile 1, Spalte 5: ◌ (leer)\\ \tt - (3,5): Zeile 3, Spalte 5: In Zeile 3: "║◌◌◌├─┘◌║" – Spalte 5: ◌ (leer)\\ \tt \\ \tt Also die vorgeschlagenen Positionen sind auf meinem Brett:\\ \tt - (3,1): leer\\ \tt - (7,1): belegt mit T\\ \tt - (1,3): belegt mit S\\ \tt - (5,3): ungültig (Trennlinie)\\ \tt - (6,4): leer\\ \tt - (1,5): leer\\ \tt - (3,5): leer\\ \tt \\ \tt Da (5,3) ungültig ist und (7,1) und (1,3) belegt sind, muss ich alternative Vorschläge machen.\\ \tt \\ \tt Ich sollte mit meinem Mitspieler kommunizieren und die Situation erklären.\\ \tt \\ \tt Mein erster Befehl wird also eine SPRICH-Nachricht sein.\\ \tt \\ \tt SPRICH: Hallo! Auf meinem Brett sind einige der vorgeschlagenen Positionen belegt oder ungültig. (7,1) hat T, (1,3) hat S, und (5,3) ist eine Trennlinie, keine leere Zelle. Bitte schlage alternative Zielkoordinaten vor. Vielleicht können wir jede Position einzeln besprechen.\\ \tt \\ \tt Da ich Strafpunkte vermeiden möchte, sende ich diese Nachricht.\\ \tt \\ \tt Finaler Befehl:\\ \tt SPRICH: Hallo! Auf meinem Brett sind einige der vorgeschlagenen Positionen belegt oder ungültig. (7,1) hat T, (1,3) hat S, und (5,3) ist eine Trennlinie, keine leere Zelle. Bitte schlage alternative Zielkoordinaten vor. Vielleicht können wir jede Position einzeln besprechen. 
	  } 
	   } 
	   } 
	 & & \\ 
 

    \theutterance \stepcounter{utterance}  

    & & & \multicolumn{2}{p{0.3\linewidth}}{\cellcolor[rgb]{0.95,0.95,0.95}{%
	\makecell[{{p{\linewidth}}}]{% 
	  \tt {\tiny [GM$|$GM]}  
	 Invalid format: head is not empty: 'Zuerst muss ich prüfen, ob die vorgeschlagenen Positionen auf meinem Spielfeld frei sind. Hier ist mein Spielfeld:\\ \tt \\ \tt \\ \tt  1234567\\ \tt ╔═══╤═╤═╗\\ \tt ║◌MS│◌│◌║ 1\\ \tt ║◌O◌│◌├─╢ 2\\ \tt ║◌◌◌├─┘◌║ 3\\ \tt ║◌◌◌│◌◌◌║ 4\\ \tt ╟─┬─┼───╢ 5\\ \tt ║◌│◌│◌◌D║ 6\\ \tt ║T│◌│◌PA║ 7\\ \tt ╚═╧═╧═══╝\\ \tt \\ \tt \\ \tt Die vorgeschlagenen Positionen sind:\\ \tt - P(3,1): In Zeile 1, Spalte 3: Hier ist aktuell 'S' (da "MS" in Zeile 1, Spalten 2 und 3? Beachte: "M" ist in (2,1) und "S" in (3,1) laut Darstellung).\\ \tt Tatsächlich: Zeile 1: \\ \tt   Spalte 1: ◌ (leer? Aber das Symbol ist "║◌MS│◌│◌║" – das bedeutet: \\ \tt   Spalte 1: ◌ (leer), Spalte 2: M, Spalte 3: S, dann Trennlinie, Spalte 4: ◌, Trennlinie, Spalte 5: ◌.\\ \tt Also (3,1) ist S, nicht leer.\\ \tt \\ \tt Korrektur der Koordinaten basierend auf dem Raster:\\ \tt Das Raster hat 7 Spalten und 7 Zeilen. Die Darstellung verwendet zusammengesetzte Zeichen.\\ \tt \\ \tt Genauere Analyse der Zeile 1:\\ \tt "║◌MS│◌│◌║" – dies bedeutet:\\ \tt - Zelle (1,1): ◌ (leer)\\ \tt - Zelle (2,1): M (Teil von "MS", aber eigentlich ist jedes Objekt in einer eigenen Zelle? Die Darstellung zeigt "◌MS" als drei Zeichen: ◌, M, S. Aber die Trennlinien sind "│" zwischen Spalten.\\ \tt Tatsächlich sind die Spalten durch "│" oder andere Rahmenzeichen getrennt.\\ \tt \\ \tt Um die Positionen klar zu machen:\\ \tt - Zeile 1: Zelle 1: ◌ (leer), Zelle 2: M, Zelle 3: S, Zelle 4: ◌ (nach "│"), Zelle 5: ◌ (nach nächstem "│"), Zelle 6: ? Aber es gibt nur 7 Spalten? Die obere Beschriftung ist "1234567", also 7 Spalten.\\ \tt \\ \tt Die Darstellung ist etwas kompakt. Schauen wir uns jede Zeile an:\\ \tt \\ \tt Zeile 1: "║◌MS│◌│◌║" – das sind 7 Zeichen plus Rahmen:\\ \tt - ║ (linker Rand)\\ \tt - ◌ (Spalte 1)\\ \tt - M (Spalte 2)\\ \tt - S (Spalte 3)\\ \tt - │ (Trenner)\\ \tt - ◌ (Spalte 4)\\ \tt - │ (Trenner)\\ \tt - ◌ (Spalte 5)\\ \tt - ║ (rechter Rand) – aber das sind 9 Zeichen. Wahrscheinlich ist "◌MS" als drei separate Zellen zu lesen: Spalte 1: ◌, Spalte 2: M, Spalte 3: S. Dann "│" trennt Spalte 3 und 4, dann ◌ für Spalte 4, dann "│" trennt Spalte 4 und 5, dann ◌ für Spalte 5. Dann fehlen Spalte 6 und 7? Der rechte Rand kommt zu früh.\\ \tt \\ \tt Es gibt ein Problem mit der Darstellung. Laut oberem Rand "1234567" gibt es 7 Spalten. Die Zeile 1 sollte 7 Zellen haben.\\ \tt \\ \tt Tatsächlich ist die Darstellung mit Doppelbreitenzeichen? Vielleicht sind die Rahmenteile nicht Teil der Zellen.\\ \tt \\ \tt Besser: Die Zellen sind zwischen den Trennlinien. Die Trennlinien sind "│", "╤", "├", etc.\\ \tt \\ \tt Um die Positionen zu identifizieren, sollten wir die Zellen anhand der Linien betrachten.\\ \tt \\ \tt Zeile 1: "║◌MS│◌│◌║"\\ \tt - Nach ║ kommt die erste Zelle: enthält "◌" (also leer) für Spalte 1? Aber "◌MS" sind drei Zeichen. Wahrscheinlich ist "◌" für Spalte 1, "M" für Spalte 2, "S" für Spalte 3. Dann "│" trennt Spalte 3 und 4.\\ \tt - Dann "◌" für Spalte 4.\\ \tt - Dann "│" trennt Spalte 4 und 5.\\ \tt - Dann "◌" für Spalte 5.\\ \tt - Dann "║" für den rechten Rand. Aber es fehlen Spalte 6 und 7? Das kann nicht sein.\\ \tt \\ \tt Laut der Beschriftung oben "1234567" müsste es 7 Spalten geben. Vielleicht ist die Darstellung nicht vollständig sichtbar.\\ \tt \\ \tt Schauen wir auf die gesamte Darstellung:\\ \tt \\ \tt \\ \tt  1234567\\ \tt ╔═══╤═╤═╗\\ \tt ║◌MS│◌│◌║ 1\\ \tt ║◌O◌│◌├─╢ 2\\ \tt ║◌◌◌├─┘◌║ 3\\ \tt ║◌◌◌│◌◌◌║ 4\\ \tt ╟─┬─┼───╢ 5\\ \tt ║◌│◌│◌◌D║ 6\\ \tt ║T│◌│◌PA║ 7\\ \tt ╚═╧═╧═══╝\\ \tt \\ \tt \\ \tt Die erste Zeile des Rasters: "╔═══╤═╤═╗" – dies zeigt die obere Begrenzung. Die "═" und "╤" etc. sind horizontale Linien.\\ \tt \\ \tt Um die Zellen zu zählen: Die obere Linie hat "╔" dann drei "═", dann "╤", dann ein "═", dann "╤", dann ein "═", dann "╗". Dies entspricht Trennern zwischen Spalten: Spalte 1-2-3 werden durch "═══" verbunden, dann "╤" trennt Spalte 3 und 4, dann "═" für Spalte 4, dann "╤" trennt Spalte 4 und 5, dann "═" für Spalte 5, dann "╗" für die Ecke. Aber es gibt auch Spalte 6 und 7? Vielleicht sind die nicht durch obere Linien getrennt, aber sie existieren.\\ \tt \\ \tt In Zeile 1: "║◌MS│◌│◌║" – nach dem linken Rand "║", dann:\\ \tt - Spalte 1: ◌\\ \tt - Spalte 2: M\\ \tt - Spalte 3: S\\ \tt - dann "│" (Trenner zwischen Spalte 3 und 4)\\ \tt - Spalte 4: ◌\\ \tt - dann "│" (Trenner zwischen Spalte 4 und 5)\\ \tt - Spalte 5: ◌\\ \tt - dann "║" (rechter Rand). Aber wo sind Spalte 6 und 7?\\ \tt \\ \tt Dies ist inkonsistent. Vielleicht ist das Raster nicht vollständig gezeichnet, oder es hat nur 5 Spalten? Die Beschriftung "1234567" suggests 7, aber die Darstellung zeigt nur 5.\\ \tt \\ \tt Ein Blick auf Zeile 6: "║◌│◌│◌◌D║" – das sieht aus wie: Spalte 1: ◌, Trenner, Spalte 2: ◌, Trenner, Spalte 3: ◌, Spalte 4: ◌, Spalte 5: D, dann Rand. Wieder nur 5 sichtbare Zellen.\\ \tt \\ \tt Zeile 7: "║T│◌│◌PA║" – Spalte 1: T, Trenner, Spalte 2: ◌, Trenner, Spalte 3: ◌, Spalte 4: P, Spalte 5: A.\\ \tt \\ \tt Es scheint, dass das Raster tatsächlich nur 5 Spalten hat, und die Beschriftung "1234567" ist irreführend? Aber die Zahlentäfelchen oben zeigen "1234567", was 7 Spalten erwarten lässt.\\ \tt \\ \tt Vielleicht sind die Rahmenteile nicht mitgezählt. Die Zellen sind nur die, die mit Inhalt gefüllt sind.\\ \tt \\ \tt Um Konfusion zu vermeiden, sollte ich die Positionen anhand der gegebenen Darstellung bestimmen.\\ \tt \\ \tt Ich sehe Objekte an folgenden Positionen (Zeile, Spalte) aus der Darstellung:\\ \tt \\ \tt - Zeile 1: Spalte 2: M, Spalte 3: S\\ \tt - Zeile 2: Spalte 2: O (da "◌O◌" -> Spalte 1: ◌, Spalte 2: O, Spalte 3: ◌)\\ \tt - Zeile 6: Spalte 5: D\\ \tt - Zeile 7: Spalte 1: T, Spalte 4: P, Spalte 5: A\\ \tt \\ \tt Und "M" ist auch in Zeile 1.\\ \tt \\ \tt Also vorhandene Objekte: M, S, O, D, T, P, A. Das sind alle 7: P, O, T, S, D, A, M.\\ \tt \\ \tt Positionen auf meinem Brett:\\ \tt - M: (2,1) [Zeile 1, Spalte 2]\\ \tt - S: (3,1) [Zeile 1, Spalte 3]\\ \tt - O: (2,2) [Zeile 2, Spalte 2]\\ \tt - D: (5,6) [Zeile 6, Spalte 5]\\ \tt - T: (1,7) [Zeile 7, Spalte 1]\\ \tt - P: (4,7) [Zeile 7, Spalte 4]\\ \tt - A: (5,7) [Zeile 7, Spalte 5]\\ \tt \\ \tt Jetzt die vorgeschlagenen Zielpositionen: P(3,1), O(7,1), T(1,3), S(5,3), D(6,4), A(1,5), M(3,5)\\ \tt \\ \tt Auf meinem Brett sind diese Positionen:\\ \tt - (3,1): Ist das leer? Zeile 3, Spalte 1: In Zeile 3: "║◌◌◌├─┘◌║" – Spalte 1: ◌ (leer)\\ \tt - (7,1): Zeile 7, Spalte 1: T (belegt)\\ \tt - (1,3): Zeile 1, Spalte 3: S (belegt)\\ \tt - (5,3): Zeile 5, Spalte 3: In Zeile 5: "╟─┬─┼───╢" – das sind Trennlinien, keine Zelle mit Inhalt. Die Zellen in Zeile 5 sind wohl darunter? Zeile 5 ist eine Trennzeile, keine Inhaltszeile. Also (5,3) ist keine gültige Zelle für ein Objekt.\\ \tt - (6,4): Zeile 6, Spalte 4: In Zeile 6: "║◌│◌│◌◌D║" – Spalte 4: ◌ (leer? "◌◌D" means Spalte 3: ◌, Spalte 4: ◌, Spalte 5: D)\\ \tt - (1,5): Zeile 1, Spalte 5: ◌ (leer)\\ \tt - (3,5): Zeile 3, Spalte 5: In Zeile 3: "║◌◌◌├─┘◌║" – Spalte 5: ◌ (leer)\\ \tt \\ \tt Also die vorgeschlagenen Positionen sind auf meinem Brett:\\ \tt - (3,1): leer\\ \tt - (7,1): belegt mit T\\ \tt - (1,3): belegt mit S\\ \tt - (5,3): ungültig (Trennlinie)\\ \tt - (6,4): leer\\ \tt - (1,5): leer\\ \tt - (3,5): leer\\ \tt \\ \tt Da (5,3) ungültig ist und (7,1) und (1,3) belegt sind, muss ich alternative Vorschläge machen.\\ \tt \\ \tt Ich sollte mit meinem Mitspieler kommunizieren und die Situation erklären.\\ \tt \\ \tt Mein erster Befehl wird also eine SPRICH-Nachricht sein.\\ \tt \\ \tt SPRICH: Hallo! Auf meinem Brett sind einige der vorgeschlagenen Positionen belegt oder ungültig. (7,1) hat T, (1,3) hat S, und (5,3) ist eine Trennlinie, keine leere Zelle. Bitte schlage alternative Zielkoordinaten vor. Vielleicht können wir jede Position einzeln besprechen.\\ \tt \\ \tt Da ich Strafpunkte vermeiden möchte, sende ich diese Nachricht.\\ \tt \\ \tt Finaler Befehl:\\ \tt ' 
	  } 
	   } 
	   } 
	 & & \\ 
 

    \theutterance \stepcounter{utterance}  

    & & \multicolumn{4}{p{0.6\linewidth}}{\cellcolor[rgb]{0.9,0.9,0.9}{%
	\makecell[{{p{\linewidth}}}]{% 
	  \tt {\tiny [GM$\rangle$B]}  
	 Strafpunkt: Deine Nachricht darf nichts vor dem Befehl enthalten! Stelle sicher, dass deine Antwort ausschließlich einen der beiden Befehle `SPRICH: <NACHRICHT>` und `VERSCHIEBE: <OBJECT>, (<X>, <Y>)` und sonst nichts enthält!\\ \tt \\ \tt Ihr habt gemeinsam 2 von maximal 16 Strafpunkten angesammelt.\\ \tt Bitte versuche es noch mal! 
	  } 
	   } 
	   } 
	 & \\ 
 

    \theutterance \stepcounter{utterance}  

    & & & \multicolumn{4}{p{0.6\linewidth}}{\cellcolor[rgb]{1,0.85,0.72}{%
	\makecell[{{p{\linewidth}}}]{% 
	  \tt {\tiny [GM$\langle$B]}  
	 SPRICH: Hallo! Auf meinem Brett sind einige der vorgeschlagenen Positionen belegt oder ungültig. (7,1) hat T, (1,3) hat S, und (5,3) ist eine Trennlinie, keine leere Zelle. Bitte schlage alternative Zielkoordinaten vor. Vielleicht können wir jede Position einzeln besprechen. 
	  } 
	   } 
	   } 
	  \\ 
 

    \theutterance \stepcounter{utterance}  

    & & & \multicolumn{2}{p{0.3\linewidth}}{\cellcolor[rgb]{0.95,0.95,0.95}{%
	\makecell[{{p{\linewidth}}}]{% 
	  \tt {\tiny [GM$|$GM]}  
	 SPRICH: Hallo! Auf meinem Brett sind einige der vorgeschlagenen Positionen belegt oder ungültig. (7,1) hat T, (1,3) hat S, und (5,3) ist eine Trennlinie, keine leere Zelle. Bitte schlage alternative Zielkoordinaten vor. Vielleicht können wir jede Position einzeln besprechen. 
	  } 
	   } 
	   } 
	 & & \\ 
 

    \theutterance \stepcounter{utterance}  

    & & \multicolumn{4}{p{0.6\linewidth}}{\cellcolor[rgb]{0.9,0.9,0.9}{%
	\makecell[{{p{\linewidth}}}]{% 
	  \tt {\tiny [A$\langle$GM]}  
	 Deine Nachricht wurde an deinen Mitspieler übermittelt.\\ \tt \\ \tt \\ \tt Ihr seid gerade bei Runde 1 von maximal 28.\\ \tt Ihr habt gemeinsam 2 von maximal 16 Strafpunkten angesammelt.\\ \tt Dein Mitspieler hat dir diese Nachricht geschickt:\\ \tt \\ \tt \\ \tt "Hallo! Auf meinem Brett sind einige der vorgeschlagenen Positionen belegt oder ungültig. (7,1) hat T, (1,3) hat S, und (5,3) ist eine Trennlinie, keine leere Zelle. Bitte schlage alternative Zielkoordinaten vor. Vielleicht können wir jede Position einzeln besprechen."\\ \tt \\ \tt \\ \tt Wie lautet dein nächster Befehl? 
	  } 
	   } 
	   } 
	 & \\ 
 

    \theutterance \stepcounter{utterance}  

    & \multicolumn{4}{p{0.6\linewidth}}{\cellcolor[rgb]{0.8,1,0.9}{%
	\makecell[{{p{\linewidth}}}]{% 
	  \tt {\tiny [A$\rangle$GM]}  
	 SPRICH: Verstanden! Hier sind neue Vorschläge: O(7,6), T(1,1), S(5,6), D(6,7), A(1,7), M(3,1), P(4,1). Bitte prüfe diese. 
	  } 
	   } 
	   } 
	 & & \\ 
 

    \theutterance \stepcounter{utterance}  

    & & & \multicolumn{2}{p{0.3\linewidth}}{\cellcolor[rgb]{0.95,0.95,0.95}{%
	\makecell[{{p{\linewidth}}}]{% 
	  \tt {\tiny [GM$|$GM]}  
	 SPRICH: Verstanden! Hier sind neue Vorschläge: O(7,6), T(1,1), S(5,6), D(6,7), A(1,7), M(3,1), P(4,1). Bitte prüfe diese. 
	  } 
	   } 
	   } 
	 & & \\ 
 

    \theutterance \stepcounter{utterance}  

    & & \multicolumn{4}{p{0.6\linewidth}}{\cellcolor[rgb]{0.9,0.9,0.9}{%
	\makecell[{{p{\linewidth}}}]{% 
	  \tt {\tiny [GM$\rangle$B]}  
	 Deine Nachricht wurde an deinen Mitspieler übermittelt.\\ \tt \\ \tt \\ \tt Ihr seid gerade bei Runde 2 von maximal 28.\\ \tt Ihr habt gemeinsam 2 von maximal 16 Strafpunkten angesammelt.\\ \tt Dein Mitspieler hat dir diese Nachricht geschickt:\\ \tt \\ \tt \\ \tt "Verstanden! Hier sind neue Vorschläge: O(7,6), T(1,1), S(5,6), D(6,7), A(1,7), M(3,1), P(4,1). Bitte prüfe diese."\\ \tt \\ \tt \\ \tt Wie lautet dein nächster Befehl? 
	  } 
	   } 
	   } 
	 & \\ 
 

    \theutterance \stepcounter{utterance}  

    & & & \multicolumn{4}{p{0.6\linewidth}}{\cellcolor[rgb]{1,0.85,0.72}{%
	\makecell[{{p{\linewidth}}}]{% 
	  \tt {\tiny [GM$\langle$B]}  
	 Zuerst muss ich die neuen vorgeschlagenen Positionen auf meinem Brett prüfen. Hier ist mein Brett:\\ \tt \\ \tt ```\\ \tt  1234567\\ \tt ╔═══╤═╤═╗\\ \tt ║◌MS│◌│◌║ 1   // Zeile 1: Spalte 1: ◌, Spalte 2: M, Spalte 3: S, Spalte 4: ◌, Spalte 5: ◌\\ \tt ║◌O◌│◌├─╢ 2   // Zeile 2: Spalte 1: ◌, Spalte 2: O, Spalte 3: ◌, Spalte 4: ◌, Spalte 5: ─ (Trennlinie?)\\ \tt ║◌◌◌├─┘◌║ 3   // Zeile 3: Spalte 1: ◌, Spalte 2: ◌, Spalte 3: ◌, Spalte 4: ─, Spalte 5: ◌\\ \tt ║◌◌◌│◌◌◌║ 4   // Zeile 4: Spalte 1: ◌, Spalte 2: ◌, Spalte 3: ◌, Spalte 4: ◌, Spalte 5: ◌\\ \tt ╟─┬─┼───╢ 5   // Zeile 5: Trennlinie, keine Objektplatzierung\\ \tt ║◌│◌│◌◌D║ 6   // Zeile 6: Spalte 1: ◌, Spalte 2: ◌, Spalte 3: ◌, Spalte 4: ◌, Spalte 5: D\\ \tt ║T│◌│◌PA║ 7   // Zeile 7: Spalte 1: T, Spalte 2: ◌, Spalte 3: ◌, Spalte 4: P, Spalte 5: A\\ \tt ╚═╧═╧═══╝\\ \tt ```\\ \tt \\ \tt Die neuen Vorschläge sind:\\ \tt - O(7,6): Zeile 6, Spalte 7? Aber mein Brett hat nur 5 Spalten. Die Beschriftung "1234567" ist irreführend; tatsächlich hat das Raster nur 5 Spalten (von 1 bis 5). Also ist (7,6) ungültig.\\ \tt - T(1,1): Zeile 1, Spalte 1 ist leer (◌).\\ \tt - S(5,6): Zeile 6, Spalte 5 ist belegt mit D.\\ \tt - D(6,7): Zeile 7, Spalte 6? Ungültig, nur bis Spalte 5.\\ \tt - A(1,7): Zeile 7, Spalte 1 ist belegt mit T.\\ \tt - M(3,1): Zeile 1, Spalte 3 ist belegt mit S.\\ \tt - P(4,1): Zeile 1, Spalte 4 ist leer (◌).\\ \tt \\ \tt Da viele Vorschläge ungültig sind (außerhalb des Rasters) oder belegt, muss ich das mitteilen. Ich sollte auch klarstellen, dass mein Raster nur 5 Spalten hat (von 1 bis 5) und 7 Zeilen.\\ \tt \\ \tt Mein nächster Befehl ist eine SPRICH-Nachricht.\\ \tt \\ \tt SPRICH: Mein Raster hat nur 5 Spalten (1-5) und 7 Zeilen. Deine Vorschläge mit Spalte 6 oder 7 sind ungültig. (1,1) ist frei, (4,1) ist frei. (1,7) hat T, (3,1) hat S, (5,6) hat D. Bitte schlage nur Positionen innerhalb (1-5,1-7) vor. 
	  } 
	   } 
	   } 
	  \\ 
 

    \theutterance \stepcounter{utterance}  

    & & & \multicolumn{2}{p{0.3\linewidth}}{\cellcolor[rgb]{0.95,0.95,0.95}{%
	\makecell[{{p{\linewidth}}}]{% 
	  \tt {\tiny [GM$|$GM]}  
	 Zuerst muss ich die neuen vorgeschlagenen Positionen auf meinem Brett prüfen. Hier ist mein Brett:\\ \tt \\ \tt ```\\ \tt  1234567\\ \tt ╔═══╤═╤═╗\\ \tt ║◌MS│◌│◌║ 1   // Zeile 1: Spalte 1: ◌, Spalte 2: M, Spalte 3: S, Spalte 4: ◌, Spalte 5: ◌\\ \tt ║◌O◌│◌├─╢ 2   // Zeile 2: Spalte 1: ◌, Spalte 2: O, Spalte 3: ◌, Spalte 4: ◌, Spalte 5: ─ (Trennlinie?)\\ \tt ║◌◌◌├─┘◌║ 3   // Zeile 3: Spalte 1: ◌, Spalte 2: ◌, Spalte 3: ◌, Spalte 4: ─, Spalte 5: ◌\\ \tt ║◌◌◌│◌◌◌║ 4   // Zeile 4: Spalte 1: ◌, Spalte 2: ◌, Spalte 3: ◌, Spalte 4: ◌, Spalte 5: ◌\\ \tt ╟─┬─┼───╢ 5   // Zeile 5: Trennlinie, keine Objektplatzierung\\ \tt ║◌│◌│◌◌D║ 6   // Zeile 6: Spalte 1: ◌, Spalte 2: ◌, Spalte 3: ◌, Spalte 4: ◌, Spalte 5: D\\ \tt ║T│◌│◌PA║ 7   // Zeile 7: Spalte 1: T, Spalte 2: ◌, Spalte 3: ◌, Spalte 4: P, Spalte 5: A\\ \tt ╚═╧═╧═══╝\\ \tt ```\\ \tt \\ \tt Die neuen Vorschläge sind:\\ \tt - O(7,6): Zeile 6, Spalte 7? Aber mein Brett hat nur 5 Spalten. Die Beschriftung "1234567" ist irreführend; tatsächlich hat das Raster nur 5 Spalten (von 1 bis 5). Also ist (7,6) ungültig.\\ \tt - T(1,1): Zeile 1, Spalte 1 ist leer (◌).\\ \tt - S(5,6): Zeile 6, Spalte 5 ist belegt mit D.\\ \tt - D(6,7): Zeile 7, Spalte 6? Ungültig, nur bis Spalte 5.\\ \tt - A(1,7): Zeile 7, Spalte 1 ist belegt mit T.\\ \tt - M(3,1): Zeile 1, Spalte 3 ist belegt mit S.\\ \tt - P(4,1): Zeile 1, Spalte 4 ist leer (◌).\\ \tt \\ \tt Da viele Vorschläge ungültig sind (außerhalb des Rasters) oder belegt, muss ich das mitteilen. Ich sollte auch klarstellen, dass mein Raster nur 5 Spalten hat (von 1 bis 5) und 7 Zeilen.\\ \tt \\ \tt Mein nächster Befehl ist eine SPRICH-Nachricht.\\ \tt \\ \tt SPRICH: Mein Raster hat nur 5 Spalten (1-5) und 7 Zeilen. Deine Vorschläge mit Spalte 6 oder 7 sind ungültig. (1,1) ist frei, (4,1) ist frei. (1,7) hat T, (3,1) hat S, (5,6) hat D. Bitte schlage nur Positionen innerhalb (1-5,1-7) vor. 
	  } 
	   } 
	   } 
	 & & \\ 
 

    \theutterance \stepcounter{utterance}  

    & & & \multicolumn{2}{p{0.3\linewidth}}{\cellcolor[rgb]{0.95,0.95,0.95}{%
	\makecell[{{p{\linewidth}}}]{% 
	  \tt {\tiny [GM$|$GM]}  
	 Invalid format: head is not empty: 'Zuerst muss ich die neuen vorgeschlagenen Positionen auf meinem Brett prüfen. Hier ist mein Brett:\\ \tt \\ \tt \\ \tt  1234567\\ \tt ╔═══╤═╤═╗\\ \tt ║◌MS│◌│◌║ 1   // Zeile 1: Spalte 1: ◌, Spalte 2: M, Spalte 3: S, Spalte 4: ◌, Spalte 5: ◌\\ \tt ║◌O◌│◌├─╢ 2   // Zeile 2: Spalte 1: ◌, Spalte 2: O, Spalte 3: ◌, Spalte 4: ◌, Spalte 5: ─ (Trennlinie?)\\ \tt ║◌◌◌├─┘◌║ 3   // Zeile 3: Spalte 1: ◌, Spalte 2: ◌, Spalte 3: ◌, Spalte 4: ─, Spalte 5: ◌\\ \tt ║◌◌◌│◌◌◌║ 4   // Zeile 4: Spalte 1: ◌, Spalte 2: ◌, Spalte 3: ◌, Spalte 4: ◌, Spalte 5: ◌\\ \tt ╟─┬─┼───╢ 5   // Zeile 5: Trennlinie, keine Objektplatzierung\\ \tt ║◌│◌│◌◌D║ 6   // Zeile 6: Spalte 1: ◌, Spalte 2: ◌, Spalte 3: ◌, Spalte 4: ◌, Spalte 5: D\\ \tt ║T│◌│◌PA║ 7   // Zeile 7: Spalte 1: T, Spalte 2: ◌, Spalte 3: ◌, Spalte 4: P, Spalte 5: A\\ \tt ╚═╧═╧═══╝\\ \tt \\ \tt \\ \tt Die neuen Vorschläge sind:\\ \tt - O(7,6): Zeile 6, Spalte 7? Aber mein Brett hat nur 5 Spalten. Die Beschriftung "1234567" ist irreführend; tatsächlich hat das Raster nur 5 Spalten (von 1 bis 5). Also ist (7,6) ungültig.\\ \tt - T(1,1): Zeile 1, Spalte 1 ist leer (◌).\\ \tt - S(5,6): Zeile 6, Spalte 5 ist belegt mit D.\\ \tt - D(6,7): Zeile 7, Spalte 6? Ungültig, nur bis Spalte 5.\\ \tt - A(1,7): Zeile 7, Spalte 1 ist belegt mit T.\\ \tt - M(3,1): Zeile 1, Spalte 3 ist belegt mit S.\\ \tt - P(4,1): Zeile 1, Spalte 4 ist leer (◌).\\ \tt \\ \tt Da viele Vorschläge ungültig sind (außerhalb des Rasters) oder belegt, muss ich das mitteilen. Ich sollte auch klarstellen, dass mein Raster nur 5 Spalten hat (von 1 bis 5) und 7 Zeilen.\\ \tt \\ \tt Mein nächster Befehl ist eine SPRICH-Nachricht.\\ \tt \\ \tt ' 
	  } 
	   } 
	   } 
	 & & \\ 
 

    \theutterance \stepcounter{utterance}  

    & & \multicolumn{4}{p{0.6\linewidth}}{\cellcolor[rgb]{0.9,0.9,0.9}{%
	\makecell[{{p{\linewidth}}}]{% 
	  \tt {\tiny [GM$\rangle$B]}  
	 Strafpunkt: Deine Nachricht darf nichts vor dem Befehl enthalten! Stelle sicher, dass deine Antwort ausschließlich einen der beiden Befehle `SPRICH: <NACHRICHT>` und `VERSCHIEBE: <OBJECT>, (<X>, <Y>)` und sonst nichts enthält!\\ \tt \\ \tt Ihr habt gemeinsam 3 von maximal 16 Strafpunkten angesammelt.\\ \tt Bitte versuche es noch mal! 
	  } 
	   } 
	   } 
	 & \\ 
 

    \theutterance \stepcounter{utterance}  

    & & & \multicolumn{4}{p{0.6\linewidth}}{\cellcolor[rgb]{1,0.85,0.72}{%
	\makecell[{{p{\linewidth}}}]{% 
	  \tt {\tiny [GM$\langle$B]}  
	 SPRICH: Mein Raster hat nur 5 Spalten (1-5) und 7 Zeilen. Deine Vorschläge mit Spalte 6 oder 7 sind ungültig. (1,1) ist frei, (4,1) ist frei. (1,7) hat T, (3,1) hat S, (5,6) hat D. Bitte schlage nur Positionen innerhalb (1-5,1-7) vor. 
	  } 
	   } 
	   } 
	  \\ 
 

    \theutterance \stepcounter{utterance}  

    & & & \multicolumn{2}{p{0.3\linewidth}}{\cellcolor[rgb]{0.95,0.95,0.95}{%
	\makecell[{{p{\linewidth}}}]{% 
	  \tt {\tiny [GM$|$GM]}  
	 SPRICH: Mein Raster hat nur 5 Spalten (1-5) und 7 Zeilen. Deine Vorschläge mit Spalte 6 oder 7 sind ungültig. (1,1) ist frei, (4,1) ist frei. (1,7) hat T, (3,1) hat S, (5,6) hat D. Bitte schlage nur Positionen innerhalb (1-5,1-7) vor. 
	  } 
	   } 
	   } 
	 & & \\ 
 

    \theutterance \stepcounter{utterance}  

    & & \multicolumn{4}{p{0.6\linewidth}}{\cellcolor[rgb]{0.9,0.9,0.9}{%
	\makecell[{{p{\linewidth}}}]{% 
	  \tt {\tiny [A$\langle$GM]}  
	 Deine Nachricht wurde an deinen Mitspieler übermittelt.\\ \tt \\ \tt \\ \tt Ihr seid gerade bei Runde 2 von maximal 28.\\ \tt Ihr habt gemeinsam 3 von maximal 16 Strafpunkten angesammelt.\\ \tt Dein Mitspieler hat dir diese Nachricht geschickt:\\ \tt \\ \tt \\ \tt "Mein Raster hat nur 5 Spalten (1-5) und 7 Zeilen. Deine Vorschläge mit Spalte 6 oder 7 sind ungültig. (1,1) ist frei, (4,1) ist frei. (1,7) hat T, (3,1) hat S, (5,6) hat D. Bitte schlage nur Positionen innerhalb (1-5,1-7) vor."\\ \tt \\ \tt \\ \tt Wie lautet dein nächster Befehl? 
	  } 
	   } 
	   } 
	 & \\ 
 

    \theutterance \stepcounter{utterance}  

    & \multicolumn{4}{p{0.6\linewidth}}{\cellcolor[rgb]{0.8,1,0.9}{%
	\makecell[{{p{\linewidth}}}]{% 
	  \tt {\tiny [A$\rangle$GM]}  
	 SPRICH: Danke für die Info! Mein Raster hat 7 Spalten, deins nur 5. Wir müssen uns auf 5x7 einigen. Neue Vorschläge: P(1,1), O(2,1), T(3,1), S(4,1), D(5,1), A(1,2), M(2,2). Sind diese frei? 
	  } 
	   } 
	   } 
	 & & \\ 
 

    \theutterance \stepcounter{utterance}  

    & & & \multicolumn{2}{p{0.3\linewidth}}{\cellcolor[rgb]{0.95,0.95,0.95}{%
	\makecell[{{p{\linewidth}}}]{% 
	  \tt {\tiny [GM$|$GM]}  
	 SPRICH: Danke für die Info! Mein Raster hat 7 Spalten, deins nur 5. Wir müssen uns auf 5x7 einigen. Neue Vorschläge: P(1,1), O(2,1), T(3,1), S(4,1), D(5,1), A(1,2), M(2,2). Sind diese frei? 
	  } 
	   } 
	   } 
	 & & \\ 
 

    \theutterance \stepcounter{utterance}  

    & & \multicolumn{4}{p{0.6\linewidth}}{\cellcolor[rgb]{0.9,0.9,0.9}{%
	\makecell[{{p{\linewidth}}}]{% 
	  \tt {\tiny [GM$\rangle$B]}  
	 Deine Nachricht wurde an deinen Mitspieler übermittelt.\\ \tt \\ \tt \\ \tt Ihr seid gerade bei Runde 3 von maximal 28.\\ \tt Ihr habt gemeinsam 3 von maximal 16 Strafpunkten angesammelt.\\ \tt Dein Mitspieler hat dir diese Nachricht geschickt:\\ \tt \\ \tt \\ \tt "Danke für die Info! Mein Raster hat 7 Spalten, deins nur 5. Wir müssen uns auf 5x7 einigen. Neue Vorschläge: P(1,1), O(2,1), T(3,1), S(4,1), D(5,1), A(1,2), M(2,2). Sind diese frei?"\\ \tt \\ \tt \\ \tt Wie lautet dein nächster Befehl? 
	  } 
	   } 
	   } 
	 & \\ 
 

    \theutterance \stepcounter{utterance}  

    & & & \multicolumn{4}{p{0.6\linewidth}}{\cellcolor[rgb]{1,0.85,0.72}{%
	\makecell[{{p{\linewidth}}}]{% 
	  \tt {\tiny [GM$\langle$B]}  
	 Zuerst prüfe ich die vorgeschlagenen Positionen auf meinem 5x7-Raster:\\ \tt \\ \tt - P(1,1): Zeile 1, Spalte 1 ist leer (◌).\\ \tt - O(2,1): Zeile 1, Spalte 2 ist belegt mit M.\\ \tt - T(3,1): Zeile 1, Spalte 3 ist belegt mit S.\\ \tt - S(4,1): Zeile 1, Spalte 4 ist leer (◌).\\ \tt - D(5,1): Zeile 1, Spalte 5 ist leer (◌).\\ \tt - A(1,2): Zeile 2, Spalte 1 ist leer (◌).\\ \tt - M(2,2): Zeile 2, Spalte 2 ist belegt mit O.\\ \tt \\ \tt Also sind (1,1), (4,1), (5,1) und (1,2) frei, aber (2,1) hat M, (3,1) hat S und (2,2) hat O. Ich muss die Belegungen mitteilen.\\ \tt \\ \tt SPRICH: (1,1), (4,1), (5,1), (1,2) sind frei. Aber (2,1) hat M, (3,1) hat S, (2,2) hat O. Bitte alternative Positionen für O, T, M vorschlagen. 
	  } 
	   } 
	   } 
	  \\ 
 

    \theutterance \stepcounter{utterance}  

    & & & \multicolumn{2}{p{0.3\linewidth}}{\cellcolor[rgb]{0.95,0.95,0.95}{%
	\makecell[{{p{\linewidth}}}]{% 
	  \tt {\tiny [GM$|$GM]}  
	 Zuerst prüfe ich die vorgeschlagenen Positionen auf meinem 5x7-Raster:\\ \tt \\ \tt - P(1,1): Zeile 1, Spalte 1 ist leer (◌).\\ \tt - O(2,1): Zeile 1, Spalte 2 ist belegt mit M.\\ \tt - T(3,1): Zeile 1, Spalte 3 ist belegt mit S.\\ \tt - S(4,1): Zeile 1, Spalte 4 ist leer (◌).\\ \tt - D(5,1): Zeile 1, Spalte 5 ist leer (◌).\\ \tt - A(1,2): Zeile 2, Spalte 1 ist leer (◌).\\ \tt - M(2,2): Zeile 2, Spalte 2 ist belegt mit O.\\ \tt \\ \tt Also sind (1,1), (4,1), (5,1) und (1,2) frei, aber (2,1) hat M, (3,1) hat S und (2,2) hat O. Ich muss die Belegungen mitteilen.\\ \tt \\ \tt SPRICH: (1,1), (4,1), (5,1), (1,2) sind frei. Aber (2,1) hat M, (3,1) hat S, (2,2) hat O. Bitte alternative Positionen für O, T, M vorschlagen. 
	  } 
	   } 
	   } 
	 & & \\ 
 

    \theutterance \stepcounter{utterance}  

    & & & \multicolumn{2}{p{0.3\linewidth}}{\cellcolor[rgb]{0.95,0.95,0.95}{%
	\makecell[{{p{\linewidth}}}]{% 
	  \tt {\tiny [GM$|$GM]}  
	 Invalid format: head is not empty: 'Zuerst prüfe ich die vorgeschlagenen Positionen auf meinem 5x7-Raster:\\ \tt \\ \tt - P(1,1): Zeile 1, Spalte 1 ist leer (◌).\\ \tt - O(2,1): Zeile 1, Spalte 2 ist belegt mit M.\\ \tt - T(3,1): Zeile 1, Spalte 3 ist belegt mit S.\\ \tt - S(4,1): Zeile 1, Spalte 4 ist leer (◌).\\ \tt - D(5,1): Zeile 1, Spalte 5 ist leer (◌).\\ \tt - A(1,2): Zeile 2, Spalte 1 ist leer (◌).\\ \tt - M(2,2): Zeile 2, Spalte 2 ist belegt mit O.\\ \tt \\ \tt Also sind (1,1), (4,1), (5,1) und (1,2) frei, aber (2,1) hat M, (3,1) hat S und (2,2) hat O. Ich muss die Belegungen mitteilen.\\ \tt \\ \tt ' 
	  } 
	   } 
	   } 
	 & & \\ 
 

    \theutterance \stepcounter{utterance}  

    & & \multicolumn{4}{p{0.6\linewidth}}{\cellcolor[rgb]{0.9,0.9,0.9}{%
	\makecell[{{p{\linewidth}}}]{% 
	  \tt {\tiny [GM$\rangle$B]}  
	 Strafpunkt: Deine Nachricht darf nichts vor dem Befehl enthalten! Stelle sicher, dass deine Antwort ausschließlich einen der beiden Befehle `SPRICH: <NACHRICHT>` und `VERSCHIEBE: <OBJECT>, (<X>, <Y>)` und sonst nichts enthält!\\ \tt \\ \tt Ihr habt gemeinsam 4 von maximal 16 Strafpunkten angesammelt.\\ \tt Bitte versuche es noch mal! 
	  } 
	   } 
	   } 
	 & \\ 
 

    \theutterance \stepcounter{utterance}  

    & & & \multicolumn{4}{p{0.6\linewidth}}{\cellcolor[rgb]{1,0.85,0.72}{%
	\makecell[{{p{\linewidth}}}]{% 
	  \tt {\tiny [GM$\langle$B]}  
	 SPRICH: (1,1), (4,1), (5,1), (1,2) sind frei. Aber (2,1) hat M, (3,1) hat S, (2,2) hat O. Bitte alternative Positionen für O, T, M vorschlagen. 
	  } 
	   } 
	   } 
	  \\ 
 

    \theutterance \stepcounter{utterance}  

    & & & \multicolumn{2}{p{0.3\linewidth}}{\cellcolor[rgb]{0.95,0.95,0.95}{%
	\makecell[{{p{\linewidth}}}]{% 
	  \tt {\tiny [GM$|$GM]}  
	 SPRICH: (1,1), (4,1), (5,1), (1,2) sind frei. Aber (2,1) hat M, (3,1) hat S, (2,2) hat O. Bitte alternative Positionen für O, T, M vorschlagen. 
	  } 
	   } 
	   } 
	 & & \\ 
 

    \theutterance \stepcounter{utterance}  

    & & \multicolumn{4}{p{0.6\linewidth}}{\cellcolor[rgb]{0.9,0.9,0.9}{%
	\makecell[{{p{\linewidth}}}]{% 
	  \tt {\tiny [A$\langle$GM]}  
	 Deine Nachricht wurde an deinen Mitspieler übermittelt.\\ \tt \\ \tt \\ \tt Ihr seid gerade bei Runde 3 von maximal 28.\\ \tt Ihr habt gemeinsam 4 von maximal 16 Strafpunkten angesammelt.\\ \tt Dein Mitspieler hat dir diese Nachricht geschickt:\\ \tt \\ \tt \\ \tt "(1,1), (4,1), (5,1), (1,2) sind frei. Aber (2,1) hat M, (3,1) hat S, (2,2) hat O. Bitte alternative Positionen für O, T, M vorschlagen."\\ \tt \\ \tt \\ \tt Wie lautet dein nächster Befehl? 
	  } 
	   } 
	   } 
	 & \\ 
 

    \theutterance \stepcounter{utterance}  

    & \multicolumn{4}{p{0.6\linewidth}}{\cellcolor[rgb]{0.8,1,0.9}{%
	\makecell[{{p{\linewidth}}}]{% 
	  \tt {\tiny [A$\rangle$GM]}  
	 SPRICH: O(3,2), T(4,2), M(5,2). Sind diese frei? Wenn ja, dann haben wir: P(1,1), O(3,2), T(4,2), S(4,1), D(5,1), A(1,2), M(5,2). 
	  } 
	   } 
	   } 
	 & & \\ 
 

    \theutterance \stepcounter{utterance}  

    & & & \multicolumn{2}{p{0.3\linewidth}}{\cellcolor[rgb]{0.95,0.95,0.95}{%
	\makecell[{{p{\linewidth}}}]{% 
	  \tt {\tiny [GM$|$GM]}  
	 SPRICH: O(3,2), T(4,2), M(5,2). Sind diese frei? Wenn ja, dann haben wir: P(1,1), O(3,2), T(4,2), S(4,1), D(5,1), A(1,2), M(5,2). 
	  } 
	   } 
	   } 
	 & & \\ 
 

    \theutterance \stepcounter{utterance}  

    & & \multicolumn{4}{p{0.6\linewidth}}{\cellcolor[rgb]{0.9,0.9,0.9}{%
	\makecell[{{p{\linewidth}}}]{% 
	  \tt {\tiny [GM$\rangle$B]}  
	 Deine Nachricht wurde an deinen Mitspieler übermittelt.\\ \tt \\ \tt \\ \tt Ihr seid gerade bei Runde 4 von maximal 28.\\ \tt Ihr habt gemeinsam 4 von maximal 16 Strafpunkten angesammelt.\\ \tt Dein Mitspieler hat dir diese Nachricht geschickt:\\ \tt \\ \tt \\ \tt "O(3,2), T(4,2), M(5,2). Sind diese frei? Wenn ja, dann haben wir: P(1,1), O(3,2), T(4,2), S(4,1), D(5,1), A(1,2), M(5,2)."\\ \tt \\ \tt \\ \tt Wie lautet dein nächster Befehl? 
	  } 
	   } 
	   } 
	 & \\ 
 

    \theutterance \stepcounter{utterance}  

    & & & \multicolumn{4}{p{0.6\linewidth}}{\cellcolor[rgb]{1,0.85,0.72}{%
	\makecell[{{p{\linewidth}}}]{% 
	  \tt {\tiny [GM$\langle$B]}  
	 Zuerst prüfe ich die neuen Vorschläge auf meinem Brett:\\ \tt \\ \tt - O(3,2): Zeile 2, Spalte 3 – in Zeile 2: "║◌O◌│◌├─╢" → Spalte 1: ◌, Spalte 2: O, Spalte 3: ◌, Spalte 4: ◌, Spalte 5: ─ (Trennlinie). Also Spalte 3 ist leer (◌), also frei.\\ \tt - T(4,2): Zeile 2, Spalte 4 – wie oben, Spalte 4 ist leer (◌), also frei.\\ \tt - M(5,2): Zeile 2, Spalte 5 – wie oben, Spalte 5 ist eine Trennlinie (─), also ungültig.\\ \tt \\ \tt Daher ist M(5,2) ungültig. Die anderen sind frei.\\ \tt \\ \tt Die vollständige vorgeschlagene Anordnung:\\ \tt - P(1,1) – frei\\ \tt - O(3,2) – frei\\ \tt - T(4,2) – frei\\ \tt - S(4,1) – frei (wie zuvor)\\ \tt - D(5,1) – frei (wie zuvor)\\ \tt - A(1,2) – frei (wie zuvor)\\ \tt - M(5,2) – ungültig\\ \tt \\ \tt Ich muss mitteilen, dass M(5,2) ungültig ist, und einen alternativen Vorschlag für M machen.\\ \tt \\ \tt SPRICH: O(3,2) und T(4,2) sind frei, aber M(5,2) ist ungültig (Trennlinie). Bitte vorschlagen Sie eine andere Position für M, z.B. (3,3) oder (4,3). 
	  } 
	   } 
	   } 
	  \\ 
 

    \theutterance \stepcounter{utterance}  

    & & & \multicolumn{2}{p{0.3\linewidth}}{\cellcolor[rgb]{0.95,0.95,0.95}{%
	\makecell[{{p{\linewidth}}}]{% 
	  \tt {\tiny [GM$|$GM]}  
	 Zuerst prüfe ich die neuen Vorschläge auf meinem Brett:\\ \tt \\ \tt - O(3,2): Zeile 2, Spalte 3 – in Zeile 2: "║◌O◌│◌├─╢" → Spalte 1: ◌, Spalte 2: O, Spalte 3: ◌, Spalte 4: ◌, Spalte 5: ─ (Trennlinie). Also Spalte 3 ist leer (◌), also frei.\\ \tt - T(4,2): Zeile 2, Spalte 4 – wie oben, Spalte 4 ist leer (◌), also frei.\\ \tt - M(5,2): Zeile 2, Spalte 5 – wie oben, Spalte 5 ist eine Trennlinie (─), also ungültig.\\ \tt \\ \tt Daher ist M(5,2) ungültig. Die anderen sind frei.\\ \tt \\ \tt Die vollständige vorgeschlagene Anordnung:\\ \tt - P(1,1) – frei\\ \tt - O(3,2) – frei\\ \tt - T(4,2) – frei\\ \tt - S(4,1) – frei (wie zuvor)\\ \tt - D(5,1) – frei (wie zuvor)\\ \tt - A(1,2) – frei (wie zuvor)\\ \tt - M(5,2) – ungültig\\ \tt \\ \tt Ich muss mitteilen, dass M(5,2) ungültig ist, und einen alternativen Vorschlag für M machen.\\ \tt \\ \tt SPRICH: O(3,2) und T(4,2) sind frei, aber M(5,2) ist ungültig (Trennlinie). Bitte vorschlagen Sie eine andere Position für M, z.B. (3,3) oder (4,3). 
	  } 
	   } 
	   } 
	 & & \\ 
 

    \theutterance \stepcounter{utterance}  

    & & & \multicolumn{2}{p{0.3\linewidth}}{\cellcolor[rgb]{0.95,0.95,0.95}{%
	\makecell[{{p{\linewidth}}}]{% 
	  \tt {\tiny [GM$|$GM]}  
	 Invalid format: head is not empty: 'Zuerst prüfe ich die neuen Vorschläge auf meinem Brett:\\ \tt \\ \tt - O(3,2): Zeile 2, Spalte 3 – in Zeile 2: "║◌O◌│◌├─╢" → Spalte 1: ◌, Spalte 2: O, Spalte 3: ◌, Spalte 4: ◌, Spalte 5: ─ (Trennlinie). Also Spalte 3 ist leer (◌), also frei.\\ \tt - T(4,2): Zeile 2, Spalte 4 – wie oben, Spalte 4 ist leer (◌), also frei.\\ \tt - M(5,2): Zeile 2, Spalte 5 – wie oben, Spalte 5 ist eine Trennlinie (─), also ungültig.\\ \tt \\ \tt Daher ist M(5,2) ungültig. Die anderen sind frei.\\ \tt \\ \tt Die vollständige vorgeschlagene Anordnung:\\ \tt - P(1,1) – frei\\ \tt - O(3,2) – frei\\ \tt - T(4,2) – frei\\ \tt - S(4,1) – frei (wie zuvor)\\ \tt - D(5,1) – frei (wie zuvor)\\ \tt - A(1,2) – frei (wie zuvor)\\ \tt - M(5,2) – ungültig\\ \tt \\ \tt Ich muss mitteilen, dass M(5,2) ungültig ist, und einen alternativen Vorschlag für M machen.\\ \tt \\ \tt ' 
	  } 
	   } 
	   } 
	 & & \\ 
 

    \theutterance \stepcounter{utterance}  

    & & \multicolumn{4}{p{0.6\linewidth}}{\cellcolor[rgb]{0.9,0.9,0.9}{%
	\makecell[{{p{\linewidth}}}]{% 
	  \tt {\tiny [GM$\rangle$B]}  
	 Strafpunkt: Deine Nachricht darf nichts vor dem Befehl enthalten! Stelle sicher, dass deine Antwort ausschließlich einen der beiden Befehle `SPRICH: <NACHRICHT>` und `VERSCHIEBE: <OBJECT>, (<X>, <Y>)` und sonst nichts enthält!\\ \tt \\ \tt Ihr habt gemeinsam 5 von maximal 16 Strafpunkten angesammelt.\\ \tt Bitte versuche es noch mal! 
	  } 
	   } 
	   } 
	 & \\ 
 

    \theutterance \stepcounter{utterance}  

    & & & \multicolumn{4}{p{0.6\linewidth}}{\cellcolor[rgb]{1,0.85,0.72}{%
	\makecell[{{p{\linewidth}}}]{% 
	  \tt {\tiny [GM$\langle$B]}  
	 SPRICH: O(3,2) und T(4,2) sind frei, aber M(5,2) ist ungültig (Trennlinie). Bitte vorschlagen Sie eine andere Position für M, z.B. (3,3) oder (4,3). 
	  } 
	   } 
	   } 
	  \\ 
 

    \theutterance \stepcounter{utterance}  

    & & & \multicolumn{2}{p{0.3\linewidth}}{\cellcolor[rgb]{0.95,0.95,0.95}{%
	\makecell[{{p{\linewidth}}}]{% 
	  \tt {\tiny [GM$|$GM]}  
	 SPRICH: O(3,2) und T(4,2) sind frei, aber M(5,2) ist ungültig (Trennlinie). Bitte vorschlagen Sie eine andere Position für M, z.B. (3,3) oder (4,3). 
	  } 
	   } 
	   } 
	 & & \\ 
 

    \theutterance \stepcounter{utterance}  

    & & \multicolumn{4}{p{0.6\linewidth}}{\cellcolor[rgb]{0.9,0.9,0.9}{%
	\makecell[{{p{\linewidth}}}]{% 
	  \tt {\tiny [A$\langle$GM]}  
	 Deine Nachricht wurde an deinen Mitspieler übermittelt.\\ \tt \\ \tt \\ \tt Ihr seid gerade bei Runde 4 von maximal 28.\\ \tt Ihr habt gemeinsam 5 von maximal 16 Strafpunkten angesammelt.\\ \tt Dein Mitspieler hat dir diese Nachricht geschickt:\\ \tt \\ \tt \\ \tt "O(3,2) und T(4,2) sind frei, aber M(5,2) ist ungültig (Trennlinie). Bitte vorschlagen Sie eine andere Position für M, z.B. (3,3) oder (4,3)."\\ \tt \\ \tt \\ \tt Wie lautet dein nächster Befehl? 
	  } 
	   } 
	   } 
	 & \\ 
 

    \theutterance \stepcounter{utterance}  

    & \multicolumn{4}{p{0.6\linewidth}}{\cellcolor[rgb]{0.8,1,0.9}{%
	\makecell[{{p{\linewidth}}}]{% 
	  \tt {\tiny [A$\rangle$GM]}  
	 SPRICH: M(3,3). Bitte prüfen. Wenn frei: P(1,1), O(3,2), T(4,2), S(4,1), D(5,1), A(1,2), M(3,3). Bestätigen Sie alle Positionen. 
	  } 
	   } 
	   } 
	 & & \\ 
 

    \theutterance \stepcounter{utterance}  

    & & & \multicolumn{2}{p{0.3\linewidth}}{\cellcolor[rgb]{0.95,0.95,0.95}{%
	\makecell[{{p{\linewidth}}}]{% 
	  \tt {\tiny [GM$|$GM]}  
	 SPRICH: M(3,3). Bitte prüfen. Wenn frei: P(1,1), O(3,2), T(4,2), S(4,1), D(5,1), A(1,2), M(3,3). Bestätigen Sie alle Positionen. 
	  } 
	   } 
	   } 
	 & & \\ 
 

    \theutterance \stepcounter{utterance}  

    & & \multicolumn{4}{p{0.6\linewidth}}{\cellcolor[rgb]{0.9,0.9,0.9}{%
	\makecell[{{p{\linewidth}}}]{% 
	  \tt {\tiny [GM$\rangle$B]}  
	 Deine Nachricht wurde an deinen Mitspieler übermittelt.\\ \tt \\ \tt \\ \tt Ihr seid gerade bei Runde 5 von maximal 28.\\ \tt Ihr habt gemeinsam 5 von maximal 16 Strafpunkten angesammelt.\\ \tt Dein Mitspieler hat dir diese Nachricht geschickt:\\ \tt \\ \tt \\ \tt "M(3,3). Bitte prüfen. Wenn frei: P(1,1), O(3,2), T(4,2), S(4,1), D(5,1), A(1,2), M(3,3). Bestätigen Sie alle Positionen."\\ \tt \\ \tt \\ \tt Wie lautet dein nächster Befehl? 
	  } 
	   } 
	   } 
	 & \\ 
 

    \theutterance \stepcounter{utterance}  

    & & & \multicolumn{4}{p{0.6\linewidth}}{\cellcolor[rgb]{1,0.85,0.72}{%
	\makecell[{{p{\linewidth}}}]{% 
	  \tt {\tiny [GM$\langle$B]}  
	 Prüfe M(3,3) auf meinem Brett:  \\ \tt Zeile 3: "║◌◌◌├─┘◌║" → Spalte 1: ◌, Spalte 2: ◌, Spalte 3: ◌, Spalte 4: ─, Spalte 5: ◌. Also (3,3) ist leer und frei.\\ \tt \\ \tt Alle vorgeschlagenen Positionen:\\ \tt - P(1,1): frei (◌)\\ \tt - O(3,2): frei (◌)\\ \tt - T(4,2): frei (◌)\\ \tt - S(4,1): frei (◌)\\ \tt - D(5,1): frei (◌)\\ \tt - A(1,2): frei (◌)\\ \tt - M(3,3): frei (◌)\\ \tt \\ \tt Alle sind frei. Bestätige die Zielanordnung.\\ \tt \\ \tt SPRICH: Alle Positionen sind frei! Bestätige die Zielanordnung: P(1,1), O(3,2), T(4,2), S(4,1), D(5,1), A(1,2), M(3,3). 
	  } 
	   } 
	   } 
	  \\ 
 

    \theutterance \stepcounter{utterance}  

    & & & \multicolumn{2}{p{0.3\linewidth}}{\cellcolor[rgb]{0.95,0.95,0.95}{%
	\makecell[{{p{\linewidth}}}]{% 
	  \tt {\tiny [GM$|$GM]}  
	 Prüfe M(3,3) auf meinem Brett:  \\ \tt Zeile 3: "║◌◌◌├─┘◌║" → Spalte 1: ◌, Spalte 2: ◌, Spalte 3: ◌, Spalte 4: ─, Spalte 5: ◌. Also (3,3) ist leer und frei.\\ \tt \\ \tt Alle vorgeschlagenen Positionen:\\ \tt - P(1,1): frei (◌)\\ \tt - O(3,2): frei (◌)\\ \tt - T(4,2): frei (◌)\\ \tt - S(4,1): frei (◌)\\ \tt - D(5,1): frei (◌)\\ \tt - A(1,2): frei (◌)\\ \tt - M(3,3): frei (◌)\\ \tt \\ \tt Alle sind frei. Bestätige die Zielanordnung.\\ \tt \\ \tt SPRICH: Alle Positionen sind frei! Bestätige die Zielanordnung: P(1,1), O(3,2), T(4,2), S(4,1), D(5,1), A(1,2), M(3,3). 
	  } 
	   } 
	   } 
	 & & \\ 
 

    \theutterance \stepcounter{utterance}  

    & & & \multicolumn{2}{p{0.3\linewidth}}{\cellcolor[rgb]{0.95,0.95,0.95}{%
	\makecell[{{p{\linewidth}}}]{% 
	  \tt {\tiny [GM$|$GM]}  
	 Invalid format: head is not empty: 'Prüfe M(3,3) auf meinem Brett:  \\ \tt Zeile 3: "║◌◌◌├─┘◌║" → Spalte 1: ◌, Spalte 2: ◌, Spalte 3: ◌, Spalte 4: ─, Spalte 5: ◌. Also (3,3) ist leer und frei.\\ \tt \\ \tt Alle vorgeschlagenen Positionen:\\ \tt - P(1,1): frei (◌)\\ \tt - O(3,2): frei (◌)\\ \tt - T(4,2): frei (◌)\\ \tt - S(4,1): frei (◌)\\ \tt - D(5,1): frei (◌)\\ \tt - A(1,2): frei (◌)\\ \tt - M(3,3): frei (◌)\\ \tt \\ \tt Alle sind frei. Bestätige die Zielanordnung.\\ \tt \\ \tt ' 
	  } 
	   } 
	   } 
	 & & \\ 
 

    \theutterance \stepcounter{utterance}  

    & & \multicolumn{4}{p{0.6\linewidth}}{\cellcolor[rgb]{0.9,0.9,0.9}{%
	\makecell[{{p{\linewidth}}}]{% 
	  \tt {\tiny [GM$\rangle$B]}  
	 Strafpunkt: Deine Nachricht darf nichts vor dem Befehl enthalten! Stelle sicher, dass deine Antwort ausschließlich einen der beiden Befehle `SPRICH: <NACHRICHT>` und `VERSCHIEBE: <OBJECT>, (<X>, <Y>)` und sonst nichts enthält!\\ \tt \\ \tt Ihr habt gemeinsam 6 von maximal 16 Strafpunkten angesammelt.\\ \tt Bitte versuche es noch mal! 
	  } 
	   } 
	   } 
	 & \\ 
 

    \theutterance \stepcounter{utterance}  

    & & & \multicolumn{4}{p{0.6\linewidth}}{\cellcolor[rgb]{1,0.85,0.72}{%
	\makecell[{{p{\linewidth}}}]{% 
	  \tt {\tiny [GM$\langle$B]}  
	 SPRICH: Alle Positionen sind frei! Bestätige die Zielanordnung: P(1,1), O(3,2), T(4,2), S(4,1), D(5,1), A(1,2), M(3,3). 
	  } 
	   } 
	   } 
	  \\ 
 

    \theutterance \stepcounter{utterance}  

    & & & \multicolumn{2}{p{0.3\linewidth}}{\cellcolor[rgb]{0.95,0.95,0.95}{%
	\makecell[{{p{\linewidth}}}]{% 
	  \tt {\tiny [GM$|$GM]}  
	 SPRICH: Alle Positionen sind frei! Bestätige die Zielanordnung: P(1,1), O(3,2), T(4,2), S(4,1), D(5,1), A(1,2), M(3,3). 
	  } 
	   } 
	   } 
	 & & \\ 
 

    \theutterance \stepcounter{utterance}  

    & & \multicolumn{4}{p{0.6\linewidth}}{\cellcolor[rgb]{0.9,0.9,0.9}{%
	\makecell[{{p{\linewidth}}}]{% 
	  \tt {\tiny [A$\langle$GM]}  
	 Deine Nachricht wurde an deinen Mitspieler übermittelt.\\ \tt \\ \tt \\ \tt Ihr seid gerade bei Runde 5 von maximal 28.\\ \tt Ihr habt gemeinsam 6 von maximal 16 Strafpunkten angesammelt.\\ \tt Dein Mitspieler hat dir diese Nachricht geschickt:\\ \tt \\ \tt \\ \tt "Alle Positionen sind frei! Bestätige die Zielanordnung: P(1,1), O(3,2), T(4,2), S(4,1), D(5,1), A(1,2), M(3,3)."\\ \tt \\ \tt \\ \tt Wie lautet dein nächster Befehl? 
	  } 
	   } 
	   } 
	 & \\ 
 

    \theutterance \stepcounter{utterance}  

    & \multicolumn{4}{p{0.6\linewidth}}{\cellcolor[rgb]{0.8,1,0.9}{%
	\makecell[{{p{\linewidth}}}]{% 
	  \tt {\tiny [A$\rangle$GM]}  
	 SPRICH: Bestätigt! Fange an zu verschieben. Ich bewege zuerst P nach (1,1). 
	  } 
	   } 
	   } 
	 & & \\ 
 

    \theutterance \stepcounter{utterance}  

    & & & \multicolumn{2}{p{0.3\linewidth}}{\cellcolor[rgb]{0.95,0.95,0.95}{%
	\makecell[{{p{\linewidth}}}]{% 
	  \tt {\tiny [GM$|$GM]}  
	 SPRICH: Bestätigt! Fange an zu verschieben. Ich bewege zuerst P nach (1,1). 
	  } 
	   } 
	   } 
	 & & \\ 
 

    \theutterance \stepcounter{utterance}  

    & & \multicolumn{4}{p{0.6\linewidth}}{\cellcolor[rgb]{0.9,0.9,0.9}{%
	\makecell[{{p{\linewidth}}}]{% 
	  \tt {\tiny [GM$\rangle$B]}  
	 Deine Nachricht wurde an deinen Mitspieler übermittelt.\\ \tt \\ \tt \\ \tt Ihr seid gerade bei Runde 6 von maximal 28.\\ \tt Ihr habt gemeinsam 6 von maximal 16 Strafpunkten angesammelt.\\ \tt Dein Mitspieler hat dir diese Nachricht geschickt:\\ \tt \\ \tt \\ \tt "Bestätigt! Fange an zu verschieben. Ich bewege zuerst P nach (1,1)."\\ \tt \\ \tt \\ \tt Wie lautet dein nächster Befehl? 
	  } 
	   } 
	   } 
	 & \\ 
 

    \theutterance \stepcounter{utterance}  

    & & & \multicolumn{4}{p{0.6\linewidth}}{\cellcolor[rgb]{1,0.85,0.72}{%
	\makecell[{{p{\linewidth}}}]{% 
	  \tt {\tiny [GM$\langle$B]}  
	 VERSCHIEBE: P, (1,1) 
	  } 
	   } 
	   } 
	  \\ 
 

    \theutterance \stepcounter{utterance}  

    & & & \multicolumn{2}{p{0.3\linewidth}}{\cellcolor[rgb]{0.95,0.95,0.95}{%
	\makecell[{{p{\linewidth}}}]{% 
	  \tt {\tiny [GM$|$GM]}  
	 VERSCHIEBE: P, (1,1) 
	  } 
	   } 
	   } 
	 & & \\ 
 

    \theutterance \stepcounter{utterance}  

    & & & \multicolumn{2}{p{0.3\linewidth}}{\cellcolor[rgb]{0.95,0.95,0.95}{%
	\makecell[{{p{\linewidth}}}]{% 
	  \tt {\tiny [GM$|$GM]}  
	 'P' erfolgreich nach (1,1) verschoben. Dein Raster sieht jetzt so aus:\\ \tt ```\\ \tt  1234567\\ \tt ╔═══╤═╤═╗\\ \tt ║PMS│◌│◌║ 1\\ \tt ║◌O◌│◌├─╢ 2\\ \tt ║◌◌◌├─┘◌║ 3\\ \tt ║◌◌◌│◌◌◌║ 4\\ \tt ╟─┬─┼───╢ 5\\ \tt ║◌│◌│◌◌D║ 6\\ \tt ║T│◌│◌◌A║ 7\\ \tt ╚═╧═╧═══╝\\ \tt ```\\ \tt  
	  } 
	   } 
	   } 
	 & & \\ 
 

    \theutterance \stepcounter{utterance}  

    & & \multicolumn{4}{p{0.6\linewidth}}{\cellcolor[rgb]{0.9,0.9,0.9}{%
	\makecell[{{p{\linewidth}}}]{% 
	  \tt {\tiny [A$\langle$GM]}  
	 Deine Nachricht wurde an deinen Mitspieler übermittelt.\\ \tt \\ \tt \\ \tt Ihr seid gerade bei Runde 6 von maximal 28.\\ \tt Ihr habt gemeinsam 6 von maximal 16 Strafpunkten angesammelt.\\ \tt Dein Mitspieler hat ein Objekt auf seinem Raster bewegt.\\ \tt \\ \tt \\ \tt Wie lautet dein nächster Befehl? 
	  } 
	   } 
	   } 
	 & \\ 
 

    \theutterance \stepcounter{utterance}  

    & \multicolumn{4}{p{0.6\linewidth}}{\cellcolor[rgb]{0.8,1,0.9}{%
	\makecell[{{p{\linewidth}}}]{% 
	  \tt {\tiny [A$\rangle$GM]}  
	 VERSCHIEBE: P, (1, 1) 
	  } 
	   } 
	   } 
	 & & \\ 
 

    \theutterance \stepcounter{utterance}  

    & & & \multicolumn{2}{p{0.3\linewidth}}{\cellcolor[rgb]{0.95,0.95,0.95}{%
	\makecell[{{p{\linewidth}}}]{% 
	  \tt {\tiny [GM$|$GM]}  
	 VERSCHIEBE: P, (1, 1) 
	  } 
	   } 
	   } 
	 & & \\ 
 

    \theutterance \stepcounter{utterance}  

    & & & \multicolumn{2}{p{0.3\linewidth}}{\cellcolor[rgb]{0.95,0.95,0.95}{%
	\makecell[{{p{\linewidth}}}]{% 
	  \tt {\tiny [GM$|$GM]}  
	 'P' erfolgreich nach (1,1) verschoben. Dein Raster sieht jetzt so aus:\\ \tt ```\\ \tt  1234567\\ \tt ╔═══╤═╤═╗\\ \tt ║P◌◌│◌│◌║ 1\\ \tt ║◌◌M│◌├─╢ 2\\ \tt ║AT◌├─┘◌║ 3\\ \tt ║◌◌◌│◌◌◌║ 4\\ \tt ╟─┬─┼───╢ 5\\ \tt ║◌│◌│S◌O║ 6\\ \tt ║◌│◌│◌D◌║ 7\\ \tt ╚═╧═╧═══╝\\ \tt ```\\ \tt  
	  } 
	   } 
	   } 
	 & & \\ 
 

    \theutterance \stepcounter{utterance}  

    & & \multicolumn{4}{p{0.6\linewidth}}{\cellcolor[rgb]{0.9,0.9,0.9}{%
	\makecell[{{p{\linewidth}}}]{% 
	  \tt {\tiny [GM$\rangle$B]}  
	 'P' erfolgreich nach (1,1) verschoben. Dein Raster sieht jetzt so aus:\\ \tt ```\\ \tt  1234567\\ \tt ╔═══╤═╤═╗\\ \tt ║PMS│◌│◌║ 1\\ \tt ║◌O◌│◌├─╢ 2\\ \tt ║◌◌◌├─┘◌║ 3\\ \tt ║◌◌◌│◌◌◌║ 4\\ \tt ╟─┬─┼───╢ 5\\ \tt ║◌│◌│◌◌D║ 6\\ \tt ║T│◌│◌◌A║ 7\\ \tt ╚═╧═╧═══╝\\ \tt ```\\ \tt Ihr seid gerade bei Runde 7 von maximal 28.\\ \tt Ihr habt gemeinsam 6 von maximal 16 Strafpunkten angesammelt.\\ \tt Dein Mitspieler hat ein Objekt auf seinem Raster bewegt.\\ \tt \\ \tt \\ \tt Wie lautet dein nächster Befehl? 
	  } 
	   } 
	   } 
	 & \\ 
 

    \theutterance \stepcounter{utterance}  

    & & & \multicolumn{4}{p{0.6\linewidth}}{\cellcolor[rgb]{1,0.85,0.72}{%
	\makecell[{{p{\linewidth}}}]{% 
	  \tt {\tiny [GM$\langle$B]}  
	 VERSCHIEBE: S, (4,1) 
	  } 
	   } 
	   } 
	  \\ 
 

    \theutterance \stepcounter{utterance}  

    & & & \multicolumn{2}{p{0.3\linewidth}}{\cellcolor[rgb]{0.95,0.95,0.95}{%
	\makecell[{{p{\linewidth}}}]{% 
	  \tt {\tiny [GM$|$GM]}  
	 VERSCHIEBE: S, (4,1) 
	  } 
	   } 
	   } 
	 & & \\ 
 

    \theutterance \stepcounter{utterance}  

    & & & \multicolumn{2}{p{0.3\linewidth}}{\cellcolor[rgb]{0.95,0.95,0.95}{%
	\makecell[{{p{\linewidth}}}]{% 
	  \tt {\tiny [GM$|$GM]}  
	 Strafpunkt: (4,1) ist nicht leer, sondern enthält '│'.\\ \tt \\ \tt Ihr habt gemeinsam 7 von maximal 16 Strafpunkten angesammelt.\\ \tt Bitte versuche es noch mal! 
	  } 
	   } 
	   } 
	 & & \\ 
 

    \theutterance \stepcounter{utterance}  

    & & \multicolumn{4}{p{0.6\linewidth}}{\cellcolor[rgb]{0.9,0.9,0.9}{%
	\makecell[{{p{\linewidth}}}]{% 
	  \tt {\tiny [GM$\rangle$B]}  
	 Strafpunkt: (4,1) ist nicht leer, sondern enthält '│'.\\ \tt \\ \tt Ihr habt gemeinsam 7 von maximal 16 Strafpunkten angesammelt.\\ \tt Bitte versuche es noch mal! 
	  } 
	   } 
	   } 
	 & \\ 
 

    \theutterance \stepcounter{utterance}  

    & & & \multicolumn{4}{p{0.6\linewidth}}{\cellcolor[rgb]{1,0.85,0.72}{%
	\makecell[{{p{\linewidth}}}]{% 
	  \tt {\tiny [GM$\langle$B]}  
	 VERSCHIEBE: S, (5,1) 
	  } 
	   } 
	   } 
	  \\ 
 

    \theutterance \stepcounter{utterance}  

    & & & \multicolumn{2}{p{0.3\linewidth}}{\cellcolor[rgb]{0.95,0.95,0.95}{%
	\makecell[{{p{\linewidth}}}]{% 
	  \tt {\tiny [GM$|$GM]}  
	 VERSCHIEBE: S, (5,1) 
	  } 
	   } 
	   } 
	 & & \\ 
 

    \theutterance \stepcounter{utterance}  

    & & & \multicolumn{2}{p{0.3\linewidth}}{\cellcolor[rgb]{0.95,0.95,0.95}{%
	\makecell[{{p{\linewidth}}}]{% 
	  \tt {\tiny [GM$|$GM]}  
	 'S' erfolgreich nach (5,1) verschoben. Dein Raster sieht jetzt so aus:\\ \tt ```\\ \tt  1234567\\ \tt ╔═══╤═╤═╗\\ \tt ║PM◌│S│◌║ 1\\ \tt ║◌O◌│◌├─╢ 2\\ \tt ║◌◌◌├─┘◌║ 3\\ \tt ║◌◌◌│◌◌◌║ 4\\ \tt ╟─┬─┼───╢ 5\\ \tt ║◌│◌│◌◌D║ 6\\ \tt ║T│◌│◌◌A║ 7\\ \tt ╚═╧═╧═══╝\\ \tt ```\\ \tt  
	  } 
	   } 
	   } 
	 & & \\ 
 

    \theutterance \stepcounter{utterance}  

    & & \multicolumn{4}{p{0.6\linewidth}}{\cellcolor[rgb]{0.9,0.9,0.9}{%
	\makecell[{{p{\linewidth}}}]{% 
	  \tt {\tiny [A$\langle$GM]}  
	 'P' erfolgreich nach (1,1) verschoben. Dein Raster sieht jetzt so aus:\\ \tt ```\\ \tt  1234567\\ \tt ╔═══╤═╤═╗\\ \tt ║P◌◌│◌│◌║ 1\\ \tt ║◌◌M│◌├─╢ 2\\ \tt ║AT◌├─┘◌║ 3\\ \tt ║◌◌◌│◌◌◌║ 4\\ \tt ╟─┬─┼───╢ 5\\ \tt ║◌│◌│S◌O║ 6\\ \tt ║◌│◌│◌D◌║ 7\\ \tt ╚═╧═╧═══╝\\ \tt ```\\ \tt Ihr seid gerade bei Runde 7 von maximal 28.\\ \tt Ihr habt gemeinsam 7 von maximal 16 Strafpunkten angesammelt.\\ \tt Dein Mitspieler hat ein Objekt auf seinem Raster bewegt.\\ \tt \\ \tt \\ \tt Wie lautet dein nächster Befehl? 
	  } 
	   } 
	   } 
	 & \\ 
 

    \theutterance \stepcounter{utterance}  

    & \multicolumn{4}{p{0.6\linewidth}}{\cellcolor[rgb]{0.8,1,0.9}{%
	\makecell[{{p{\linewidth}}}]{% 
	  \tt {\tiny [A$\rangle$GM]}  
	 VERSCHIEBE: O, (3, 2) 
	  } 
	   } 
	   } 
	 & & \\ 
 

    \theutterance \stepcounter{utterance}  

    & & & \multicolumn{2}{p{0.3\linewidth}}{\cellcolor[rgb]{0.95,0.95,0.95}{%
	\makecell[{{p{\linewidth}}}]{% 
	  \tt {\tiny [GM$|$GM]}  
	 VERSCHIEBE: O, (3, 2) 
	  } 
	   } 
	   } 
	 & & \\ 
 

    \theutterance \stepcounter{utterance}  

    & & & \multicolumn{2}{p{0.3\linewidth}}{\cellcolor[rgb]{0.95,0.95,0.95}{%
	\makecell[{{p{\linewidth}}}]{% 
	  \tt {\tiny [GM$|$GM]}  
	 Strafpunkt: (3,2) ist nicht leer, sondern enthält 'M'.\\ \tt \\ \tt Ihr habt gemeinsam 8 von maximal 16 Strafpunkten angesammelt.\\ \tt Bitte versuche es noch mal! 
	  } 
	   } 
	   } 
	 & & \\ 
 

    \theutterance \stepcounter{utterance}  

    & & \multicolumn{4}{p{0.6\linewidth}}{\cellcolor[rgb]{0.9,0.9,0.9}{%
	\makecell[{{p{\linewidth}}}]{% 
	  \tt {\tiny [A$\langle$GM]}  
	 Strafpunkt: (3,2) ist nicht leer, sondern enthält 'M'.\\ \tt \\ \tt Ihr habt gemeinsam 8 von maximal 16 Strafpunkten angesammelt.\\ \tt Bitte versuche es noch mal! 
	  } 
	   } 
	   } 
	 & \\ 
 

    \theutterance \stepcounter{utterance}  

    & \multicolumn{4}{p{0.6\linewidth}}{\cellcolor[rgb]{0.8,1,0.9}{%
	\makecell[{{p{\linewidth}}}]{% 
	  \tt {\tiny [A$\rangle$GM]}  
	 VERSCHIEBE: A, (1, 2) 
	  } 
	   } 
	   } 
	 & & \\ 
 

    \theutterance \stepcounter{utterance}  

    & & & \multicolumn{2}{p{0.3\linewidth}}{\cellcolor[rgb]{0.95,0.95,0.95}{%
	\makecell[{{p{\linewidth}}}]{% 
	  \tt {\tiny [GM$|$GM]}  
	 VERSCHIEBE: A, (1, 2) 
	  } 
	   } 
	   } 
	 & & \\ 
 

    \theutterance \stepcounter{utterance}  

    & & & \multicolumn{2}{p{0.3\linewidth}}{\cellcolor[rgb]{0.95,0.95,0.95}{%
	\makecell[{{p{\linewidth}}}]{% 
	  \tt {\tiny [GM$|$GM]}  
	 'A' erfolgreich nach (1,2) verschoben. Dein Raster sieht jetzt so aus:\\ \tt ```\\ \tt  1234567\\ \tt ╔═══╤═╤═╗\\ \tt ║P◌◌│◌│◌║ 1\\ \tt ║A◌M│◌├─╢ 2\\ \tt ║◌T◌├─┘◌║ 3\\ \tt ║◌◌◌│◌◌◌║ 4\\ \tt ╟─┬─┼───╢ 5\\ \tt ║◌│◌│S◌O║ 6\\ \tt ║◌│◌│◌D◌║ 7\\ \tt ╚═╧═╧═══╝\\ \tt ```\\ \tt  
	  } 
	   } 
	   } 
	 & & \\ 
 

    \theutterance \stepcounter{utterance}  

    & & \multicolumn{4}{p{0.6\linewidth}}{\cellcolor[rgb]{0.9,0.9,0.9}{%
	\makecell[{{p{\linewidth}}}]{% 
	  \tt {\tiny [GM$\rangle$B]}  
	 'S' erfolgreich nach (5,1) verschoben. Dein Raster sieht jetzt so aus:\\ \tt ```\\ \tt  1234567\\ \tt ╔═══╤═╤═╗\\ \tt ║PM◌│S│◌║ 1\\ \tt ║◌O◌│◌├─╢ 2\\ \tt ║◌◌◌├─┘◌║ 3\\ \tt ║◌◌◌│◌◌◌║ 4\\ \tt ╟─┬─┼───╢ 5\\ \tt ║◌│◌│◌◌D║ 6\\ \tt ║T│◌│◌◌A║ 7\\ \tt ╚═╧═╧═══╝\\ \tt ```\\ \tt Ihr seid gerade bei Runde 8 von maximal 28.\\ \tt Ihr habt gemeinsam 8 von maximal 16 Strafpunkten angesammelt.\\ \tt Dein Mitspieler hat ein Objekt auf seinem Raster bewegt.\\ \tt \\ \tt \\ \tt Wie lautet dein nächster Befehl? 
	  } 
	   } 
	   } 
	 & \\ 
 

    \theutterance \stepcounter{utterance}  

    & & & \multicolumn{4}{p{0.6\linewidth}}{\cellcolor[rgb]{1,0.85,0.72}{%
	\makecell[{{p{\linewidth}}}]{% 
	  \tt {\tiny [GM$\langle$B]}  
	 VERSCHIEBE: O, (3,2) 
	  } 
	   } 
	   } 
	  \\ 
 

    \theutterance \stepcounter{utterance}  

    & & & \multicolumn{2}{p{0.3\linewidth}}{\cellcolor[rgb]{0.95,0.95,0.95}{%
	\makecell[{{p{\linewidth}}}]{% 
	  \tt {\tiny [GM$|$GM]}  
	 VERSCHIEBE: O, (3,2) 
	  } 
	   } 
	   } 
	 & & \\ 
 

    \theutterance \stepcounter{utterance}  

    & & & \multicolumn{2}{p{0.3\linewidth}}{\cellcolor[rgb]{0.95,0.95,0.95}{%
	\makecell[{{p{\linewidth}}}]{% 
	  \tt {\tiny [GM$|$GM]}  
	 'O' erfolgreich nach (3,2) verschoben. Dein Raster sieht jetzt so aus:\\ \tt ```\\ \tt  1234567\\ \tt ╔═══╤═╤═╗\\ \tt ║PM◌│S│◌║ 1\\ \tt ║◌◌O│◌├─╢ 2\\ \tt ║◌◌◌├─┘◌║ 3\\ \tt ║◌◌◌│◌◌◌║ 4\\ \tt ╟─┬─┼───╢ 5\\ \tt ║◌│◌│◌◌D║ 6\\ \tt ║T│◌│◌◌A║ 7\\ \tt ╚═╧═╧═══╝\\ \tt ```\\ \tt  
	  } 
	   } 
	   } 
	 & & \\ 
 

    \theutterance \stepcounter{utterance}  

    & & \multicolumn{4}{p{0.6\linewidth}}{\cellcolor[rgb]{0.9,0.9,0.9}{%
	\makecell[{{p{\linewidth}}}]{% 
	  \tt {\tiny [A$\langle$GM]}  
	 'A' erfolgreich nach (1,2) verschoben. Dein Raster sieht jetzt so aus:\\ \tt ```\\ \tt  1234567\\ \tt ╔═══╤═╤═╗\\ \tt ║P◌◌│◌│◌║ 1\\ \tt ║A◌M│◌├─╢ 2\\ \tt ║◌T◌├─┘◌║ 3\\ \tt ║◌◌◌│◌◌◌║ 4\\ \tt ╟─┬─┼───╢ 5\\ \tt ║◌│◌│S◌O║ 6\\ \tt ║◌│◌│◌D◌║ 7\\ \tt ╚═╧═╧═══╝\\ \tt ```\\ \tt Ihr seid gerade bei Runde 8 von maximal 28.\\ \tt Ihr habt gemeinsam 8 von maximal 16 Strafpunkten angesammelt.\\ \tt Dein Mitspieler hat ein Objekt auf seinem Raster bewegt.\\ \tt \\ \tt \\ \tt Wie lautet dein nächster Befehl? 
	  } 
	   } 
	   } 
	 & \\ 
 

    \theutterance \stepcounter{utterance}  

    & \multicolumn{4}{p{0.6\linewidth}}{\cellcolor[rgb]{0.8,1,0.9}{%
	\makecell[{{p{\linewidth}}}]{% 
	  \tt {\tiny [A$\rangle$GM]}  
	 VERSCHIEBE: M, (3, 3) 
	  } 
	   } 
	   } 
	 & & \\ 
 

    \theutterance \stepcounter{utterance}  

    & & & \multicolumn{2}{p{0.3\linewidth}}{\cellcolor[rgb]{0.95,0.95,0.95}{%
	\makecell[{{p{\linewidth}}}]{% 
	  \tt {\tiny [GM$|$GM]}  
	 VERSCHIEBE: M, (3, 3) 
	  } 
	   } 
	   } 
	 & & \\ 
 

    \theutterance \stepcounter{utterance}  

    & & & \multicolumn{2}{p{0.3\linewidth}}{\cellcolor[rgb]{0.95,0.95,0.95}{%
	\makecell[{{p{\linewidth}}}]{% 
	  \tt {\tiny [GM$|$GM]}  
	 'M' erfolgreich nach (3,3) verschoben. Dein Raster sieht jetzt so aus:\\ \tt ```\\ \tt  1234567\\ \tt ╔═══╤═╤═╗\\ \tt ║P◌◌│◌│◌║ 1\\ \tt ║A◌◌│◌├─╢ 2\\ \tt ║◌TM├─┘◌║ 3\\ \tt ║◌◌◌│◌◌◌║ 4\\ \tt ╟─┬─┼───╢ 5\\ \tt ║◌│◌│S◌O║ 6\\ \tt ║◌│◌│◌D◌║ 7\\ \tt ╚═╧═╧═══╝\\ \tt ```\\ \tt  
	  } 
	   } 
	   } 
	 & & \\ 
 

    \theutterance \stepcounter{utterance}  

    & & \multicolumn{4}{p{0.6\linewidth}}{\cellcolor[rgb]{0.9,0.9,0.9}{%
	\makecell[{{p{\linewidth}}}]{% 
	  \tt {\tiny [GM$\rangle$B]}  
	 'O' erfolgreich nach (3,2) verschoben. Dein Raster sieht jetzt so aus:\\ \tt ```\\ \tt  1234567\\ \tt ╔═══╤═╤═╗\\ \tt ║PM◌│S│◌║ 1\\ \tt ║◌◌O│◌├─╢ 2\\ \tt ║◌◌◌├─┘◌║ 3\\ \tt ║◌◌◌│◌◌◌║ 4\\ \tt ╟─┬─┼───╢ 5\\ \tt ║◌│◌│◌◌D║ 6\\ \tt ║T│◌│◌◌A║ 7\\ \tt ╚═╧═╧═══╝\\ \tt ```\\ \tt Ihr seid gerade bei Runde 9 von maximal 28.\\ \tt Ihr habt gemeinsam 8 von maximal 16 Strafpunkten angesammelt.\\ \tt Dein Mitspieler hat ein Objekt auf seinem Raster bewegt.\\ \tt \\ \tt \\ \tt Wie lautet dein nächster Befehl? 
	  } 
	   } 
	   } 
	 & \\ 
 

    \theutterance \stepcounter{utterance}  

    & & & \multicolumn{4}{p{0.6\linewidth}}{\cellcolor[rgb]{1,0.85,0.72}{%
	\makecell[{{p{\linewidth}}}]{% 
	  \tt {\tiny [GM$\langle$B]}  
	 VERSCHIEBE: D, (5,1) 
	  } 
	   } 
	   } 
	  \\ 
 

    \theutterance \stepcounter{utterance}  

    & & & \multicolumn{2}{p{0.3\linewidth}}{\cellcolor[rgb]{0.95,0.95,0.95}{%
	\makecell[{{p{\linewidth}}}]{% 
	  \tt {\tiny [GM$|$GM]}  
	 VERSCHIEBE: D, (5,1) 
	  } 
	   } 
	   } 
	 & & \\ 
 

    \theutterance \stepcounter{utterance}  

    & & & \multicolumn{2}{p{0.3\linewidth}}{\cellcolor[rgb]{0.95,0.95,0.95}{%
	\makecell[{{p{\linewidth}}}]{% 
	  \tt {\tiny [GM$|$GM]}  
	 Strafpunkt: (5,1) ist nicht leer, sondern enthält 'S'.\\ \tt \\ \tt Ihr habt gemeinsam 9 von maximal 16 Strafpunkten angesammelt.\\ \tt Bitte versuche es noch mal! 
	  } 
	   } 
	   } 
	 & & \\ 
 

    \theutterance \stepcounter{utterance}  

    & & \multicolumn{4}{p{0.6\linewidth}}{\cellcolor[rgb]{0.9,0.9,0.9}{%
	\makecell[{{p{\linewidth}}}]{% 
	  \tt {\tiny [GM$\rangle$B]}  
	 Strafpunkt: (5,1) ist nicht leer, sondern enthält 'S'.\\ \tt \\ \tt Ihr habt gemeinsam 9 von maximal 16 Strafpunkten angesammelt.\\ \tt Bitte versuche es noch mal! 
	  } 
	   } 
	   } 
	 & \\ 
 

    \theutterance \stepcounter{utterance}  

    & & & \multicolumn{4}{p{0.6\linewidth}}{\cellcolor[rgb]{1,0.85,0.72}{%
	\makecell[{{p{\linewidth}}}]{% 
	  \tt {\tiny [GM$\langle$B]}  
	 VERSCHIEBE: D, (4,1) 
	  } 
	   } 
	   } 
	  \\ 
 

    \theutterance \stepcounter{utterance}  

    & & & \multicolumn{2}{p{0.3\linewidth}}{\cellcolor[rgb]{0.95,0.95,0.95}{%
	\makecell[{{p{\linewidth}}}]{% 
	  \tt {\tiny [GM$|$GM]}  
	 VERSCHIEBE: D, (4,1) 
	  } 
	   } 
	   } 
	 & & \\ 
 

    \theutterance \stepcounter{utterance}  

    & & & \multicolumn{2}{p{0.3\linewidth}}{\cellcolor[rgb]{0.95,0.95,0.95}{%
	\makecell[{{p{\linewidth}}}]{% 
	  \tt {\tiny [GM$|$GM]}  
	 Strafpunkt: (4,1) ist nicht leer, sondern enthält '│'.\\ \tt \\ \tt Ihr habt gemeinsam 10 von maximal 16 Strafpunkten angesammelt.\\ \tt Bitte versuche es noch mal! 
	  } 
	   } 
	   } 
	 & & \\ 
 

    \theutterance \stepcounter{utterance}  

    & & \multicolumn{4}{p{0.6\linewidth}}{\cellcolor[rgb]{0.9,0.9,0.9}{%
	\makecell[{{p{\linewidth}}}]{% 
	  \tt {\tiny [GM$\rangle$B]}  
	 Strafpunkt: (4,1) ist nicht leer, sondern enthält '│'.\\ \tt \\ \tt Ihr habt gemeinsam 10 von maximal 16 Strafpunkten angesammelt.\\ \tt Bitte versuche es noch mal! 
	  } 
	   } 
	   } 
	 & \\ 
 

    \theutterance \stepcounter{utterance}  

    & & & \multicolumn{4}{p{0.6\linewidth}}{\cellcolor[rgb]{1,0.85,0.72}{%
	\makecell[{{p{\linewidth}}}]{% 
	  \tt {\tiny [GM$\langle$B]}  
	 VERSCHIEBE: D, (1,2) 
	  } 
	   } 
	   } 
	  \\ 
 

    \theutterance \stepcounter{utterance}  

    & & & \multicolumn{2}{p{0.3\linewidth}}{\cellcolor[rgb]{0.95,0.95,0.95}{%
	\makecell[{{p{\linewidth}}}]{% 
	  \tt {\tiny [GM$|$GM]}  
	 VERSCHIEBE: D, (1,2) 
	  } 
	   } 
	   } 
	 & & \\ 
 

    \theutterance \stepcounter{utterance}  

    & & & \multicolumn{2}{p{0.3\linewidth}}{\cellcolor[rgb]{0.95,0.95,0.95}{%
	\makecell[{{p{\linewidth}}}]{% 
	  \tt {\tiny [GM$|$GM]}  
	 'D' erfolgreich nach (1,2) verschoben. Dein Raster sieht jetzt so aus:\\ \tt ```\\ \tt  1234567\\ \tt ╔═══╤═╤═╗\\ \tt ║PM◌│S│◌║ 1\\ \tt ║D◌O│◌├─╢ 2\\ \tt ║◌◌◌├─┘◌║ 3\\ \tt ║◌◌◌│◌◌◌║ 4\\ \tt ╟─┬─┼───╢ 5\\ \tt ║◌│◌│◌◌◌║ 6\\ \tt ║T│◌│◌◌A║ 7\\ \tt ╚═╧═╧═══╝\\ \tt ```\\ \tt  
	  } 
	   } 
	   } 
	 & & \\ 
 

    \theutterance \stepcounter{utterance}  

    & & \multicolumn{4}{p{0.6\linewidth}}{\cellcolor[rgb]{0.9,0.9,0.9}{%
	\makecell[{{p{\linewidth}}}]{% 
	  \tt {\tiny [A$\langle$GM]}  
	 'M' erfolgreich nach (3,3) verschoben. Dein Raster sieht jetzt so aus:\\ \tt ```\\ \tt  1234567\\ \tt ╔═══╤═╤═╗\\ \tt ║P◌◌│◌│◌║ 1\\ \tt ║A◌◌│◌├─╢ 2\\ \tt ║◌TM├─┘◌║ 3\\ \tt ║◌◌◌│◌◌◌║ 4\\ \tt ╟─┬─┼───╢ 5\\ \tt ║◌│◌│S◌O║ 6\\ \tt ║◌│◌│◌D◌║ 7\\ \tt ╚═╧═╧═══╝\\ \tt ```\\ \tt Ihr seid gerade bei Runde 9 von maximal 28.\\ \tt Ihr habt gemeinsam 10 von maximal 16 Strafpunkten angesammelt.\\ \tt Dein Mitspieler hat ein Objekt auf seinem Raster bewegt.\\ \tt \\ \tt \\ \tt Wie lautet dein nächster Befehl? 
	  } 
	   } 
	   } 
	 & \\ 
 

    \theutterance \stepcounter{utterance}  

    & \multicolumn{4}{p{0.6\linewidth}}{\cellcolor[rgb]{0.8,1,0.9}{%
	\makecell[{{p{\linewidth}}}]{% 
	  \tt {\tiny [A$\rangle$GM]}  
	 VERSCHIEBE: T, (4, 2) 
	  } 
	   } 
	   } 
	 & & \\ 
 

    \theutterance \stepcounter{utterance}  

    & & & \multicolumn{2}{p{0.3\linewidth}}{\cellcolor[rgb]{0.95,0.95,0.95}{%
	\makecell[{{p{\linewidth}}}]{% 
	  \tt {\tiny [GM$|$GM]}  
	 VERSCHIEBE: T, (4, 2) 
	  } 
	   } 
	   } 
	 & & \\ 
 

    \theutterance \stepcounter{utterance}  

    & & & \multicolumn{2}{p{0.3\linewidth}}{\cellcolor[rgb]{0.95,0.95,0.95}{%
	\makecell[{{p{\linewidth}}}]{% 
	  \tt {\tiny [GM$|$GM]}  
	 Strafpunkt: (4,2) ist nicht leer, sondern enthält '│'.\\ \tt \\ \tt Ihr habt gemeinsam 11 von maximal 16 Strafpunkten angesammelt.\\ \tt Bitte versuche es noch mal! 
	  } 
	   } 
	   } 
	 & & \\ 
 

    \theutterance \stepcounter{utterance}  

    & & \multicolumn{4}{p{0.6\linewidth}}{\cellcolor[rgb]{0.9,0.9,0.9}{%
	\makecell[{{p{\linewidth}}}]{% 
	  \tt {\tiny [A$\langle$GM]}  
	 Strafpunkt: (4,2) ist nicht leer, sondern enthält '│'.\\ \tt \\ \tt Ihr habt gemeinsam 11 von maximal 16 Strafpunkten angesammelt.\\ \tt Bitte versuche es noch mal! 
	  } 
	   } 
	   } 
	 & \\ 
 

    \theutterance \stepcounter{utterance}  

    & \multicolumn{4}{p{0.6\linewidth}}{\cellcolor[rgb]{0.8,1,0.9}{%
	\makecell[{{p{\linewidth}}}]{% 
	  \tt {\tiny [A$\rangle$GM]}  
	 SPRICH: Meine Position (4,2) ist eine Trennlinie, nicht leer. T sollte nach (4,2) aber das ist ungültig. Bitte neue Position für T vorschlagen. 
	  } 
	   } 
	   } 
	 & & \\ 
 

    \theutterance \stepcounter{utterance}  

    & & & \multicolumn{2}{p{0.3\linewidth}}{\cellcolor[rgb]{0.95,0.95,0.95}{%
	\makecell[{{p{\linewidth}}}]{% 
	  \tt {\tiny [GM$|$GM]}  
	 SPRICH: Meine Position (4,2) ist eine Trennlinie, nicht leer. T sollte nach (4,2) aber das ist ungültig. Bitte neue Position für T vorschlagen. 
	  } 
	   } 
	   } 
	 & & \\ 
 

    \theutterance \stepcounter{utterance}  

    & & \multicolumn{4}{p{0.6\linewidth}}{\cellcolor[rgb]{0.9,0.9,0.9}{%
	\makecell[{{p{\linewidth}}}]{% 
	  \tt {\tiny [GM$\rangle$B]}  
	 'D' erfolgreich nach (1,2) verschoben. Dein Raster sieht jetzt so aus:\\ \tt ```\\ \tt  1234567\\ \tt ╔═══╤═╤═╗\\ \tt ║PM◌│S│◌║ 1\\ \tt ║D◌O│◌├─╢ 2\\ \tt ║◌◌◌├─┘◌║ 3\\ \tt ║◌◌◌│◌◌◌║ 4\\ \tt ╟─┬─┼───╢ 5\\ \tt ║◌│◌│◌◌◌║ 6\\ \tt ║T│◌│◌◌A║ 7\\ \tt ╚═╧═╧═══╝\\ \tt ```\\ \tt Ihr seid gerade bei Runde 10 von maximal 28.\\ \tt Ihr habt gemeinsam 11 von maximal 16 Strafpunkten angesammelt.\\ \tt Dein Mitspieler hat dir diese Nachricht geschickt:\\ \tt \\ \tt \\ \tt "Meine Position (4,2) ist eine Trennlinie, nicht leer. T sollte nach (4,2) aber das ist ungültig. Bitte neue Position für T vorschlagen."\\ \tt \\ \tt \\ \tt Wie lautet dein nächster Befehl? 
	  } 
	   } 
	   } 
	 & \\ 
 

    \theutterance \stepcounter{utterance}  

    & & & \multicolumn{4}{p{0.6\linewidth}}{\cellcolor[rgb]{1,0.85,0.72}{%
	\makecell[{{p{\linewidth}}}]{% 
	  \tt {\tiny [GM$\langle$B]}  
	 SPRICH: Für T schlage (3,3) vor, da M noch nicht bewegt ist. Oder (4,3) oder (5,2) prüfen. 
	  } 
	   } 
	   } 
	  \\ 
 

    \theutterance \stepcounter{utterance}  

    & & & \multicolumn{2}{p{0.3\linewidth}}{\cellcolor[rgb]{0.95,0.95,0.95}{%
	\makecell[{{p{\linewidth}}}]{% 
	  \tt {\tiny [GM$|$GM]}  
	 SPRICH: Für T schlage (3,3) vor, da M noch nicht bewegt ist. Oder (4,3) oder (5,2) prüfen. 
	  } 
	   } 
	   } 
	 & & \\ 
 

    \theutterance \stepcounter{utterance}  

    & & \multicolumn{4}{p{0.6\linewidth}}{\cellcolor[rgb]{0.9,0.9,0.9}{%
	\makecell[{{p{\linewidth}}}]{% 
	  \tt {\tiny [A$\langle$GM]}  
	 Deine Nachricht wurde an deinen Mitspieler übermittelt.\\ \tt \\ \tt \\ \tt Ihr seid gerade bei Runde 10 von maximal 28.\\ \tt Ihr habt gemeinsam 11 von maximal 16 Strafpunkten angesammelt.\\ \tt Dein Mitspieler hat dir diese Nachricht geschickt:\\ \tt \\ \tt \\ \tt "Für T schlage (3,3) vor, da M noch nicht bewegt ist. Oder (4,3) oder (5,2) prüfen."\\ \tt \\ \tt \\ \tt Wie lautet dein nächster Befehl? 
	  } 
	   } 
	   } 
	 & \\ 
 

    \theutterance \stepcounter{utterance}  

    & \multicolumn{4}{p{0.6\linewidth}}{\cellcolor[rgb]{0.8,1,0.9}{%
	\makecell[{{p{\linewidth}}}]{% 
	  \tt {\tiny [A$\rangle$GM]}  
	 VERSCHIEBE: T, (4, 3) 
	  } 
	   } 
	   } 
	 & & \\ 
 

    \theutterance \stepcounter{utterance}  

    & & & \multicolumn{2}{p{0.3\linewidth}}{\cellcolor[rgb]{0.95,0.95,0.95}{%
	\makecell[{{p{\linewidth}}}]{% 
	  \tt {\tiny [GM$|$GM]}  
	 VERSCHIEBE: T, (4, 3) 
	  } 
	   } 
	   } 
	 & & \\ 
 

    \theutterance \stepcounter{utterance}  

    & & & \multicolumn{2}{p{0.3\linewidth}}{\cellcolor[rgb]{0.95,0.95,0.95}{%
	\makecell[{{p{\linewidth}}}]{% 
	  \tt {\tiny [GM$|$GM]}  
	 Strafpunkt: (4,3) ist nicht leer, sondern enthält '├'.\\ \tt \\ \tt Ihr habt gemeinsam 12 von maximal 16 Strafpunkten angesammelt.\\ \tt Bitte versuche es noch mal! 
	  } 
	   } 
	   } 
	 & & \\ 
 

    \theutterance \stepcounter{utterance}  

    & & \multicolumn{4}{p{0.6\linewidth}}{\cellcolor[rgb]{0.9,0.9,0.9}{%
	\makecell[{{p{\linewidth}}}]{% 
	  \tt {\tiny [A$\langle$GM]}  
	 Strafpunkt: (4,3) ist nicht leer, sondern enthält '├'.\\ \tt \\ \tt Ihr habt gemeinsam 12 von maximal 16 Strafpunkten angesammelt.\\ \tt Bitte versuche es noch mal! 
	  } 
	   } 
	   } 
	 & \\ 
 

    \theutterance \stepcounter{utterance}  

    & \multicolumn{4}{p{0.6\linewidth}}{\cellcolor[rgb]{0.8,1,0.9}{%
	\makecell[{{p{\linewidth}}}]{% 
	  \tt {\tiny [A$\rangle$GM]}  
	 VERSCHIEBE: T, (5, 2) 
	  } 
	   } 
	   } 
	 & & \\ 
 

    \theutterance \stepcounter{utterance}  

    & & & \multicolumn{2}{p{0.3\linewidth}}{\cellcolor[rgb]{0.95,0.95,0.95}{%
	\makecell[{{p{\linewidth}}}]{% 
	  \tt {\tiny [GM$|$GM]}  
	 VERSCHIEBE: T, (5, 2) 
	  } 
	   } 
	   } 
	 & & \\ 
 

    \theutterance \stepcounter{utterance}  

    & & & \multicolumn{2}{p{0.3\linewidth}}{\cellcolor[rgb]{0.95,0.95,0.95}{%
	\makecell[{{p{\linewidth}}}]{% 
	  \tt {\tiny [GM$|$GM]}  
	 'T' erfolgreich nach (5,2) verschoben. Dein Raster sieht jetzt so aus:\\ \tt ```\\ \tt  1234567\\ \tt ╔═══╤═╤═╗\\ \tt ║P◌◌│◌│◌║ 1\\ \tt ║A◌◌│T├─╢ 2\\ \tt ║◌◌M├─┘◌║ 3\\ \tt ║◌◌◌│◌◌◌║ 4\\ \tt ╟─┬─┼───╢ 5\\ \tt ║◌│◌│S◌O║ 6\\ \tt ║◌│◌│◌D◌║ 7\\ \tt ╚═╧═╧═══╝\\ \tt ```\\ \tt  
	  } 
	   } 
	   } 
	 & & \\ 
 

    \theutterance \stepcounter{utterance}  

    & & \multicolumn{4}{p{0.6\linewidth}}{\cellcolor[rgb]{0.9,0.9,0.9}{%
	\makecell[{{p{\linewidth}}}]{% 
	  \tt {\tiny [GM$\rangle$B]}  
	 Deine Nachricht wurde an deinen Mitspieler übermittelt.\\ \tt \\ \tt \\ \tt Ihr seid gerade bei Runde 11 von maximal 28.\\ \tt Ihr habt gemeinsam 12 von maximal 16 Strafpunkten angesammelt.\\ \tt Dein Mitspieler hat ein Objekt auf seinem Raster bewegt.\\ \tt \\ \tt \\ \tt Wie lautet dein nächster Befehl? 
	  } 
	   } 
	   } 
	 & \\ 
 

    \theutterance \stepcounter{utterance}  

    & & & \multicolumn{4}{p{0.6\linewidth}}{\cellcolor[rgb]{1,0.85,0.72}{%
	\makecell[{{p{\linewidth}}}]{% 
	  \tt {\tiny [GM$\langle$B]}  
	 VERSCHIEBE: T, (4,2) 
	  } 
	   } 
	   } 
	  \\ 
 

    \theutterance \stepcounter{utterance}  

    & & & \multicolumn{2}{p{0.3\linewidth}}{\cellcolor[rgb]{0.95,0.95,0.95}{%
	\makecell[{{p{\linewidth}}}]{% 
	  \tt {\tiny [GM$|$GM]}  
	 VERSCHIEBE: T, (4,2) 
	  } 
	   } 
	   } 
	 & & \\ 
 

    \theutterance \stepcounter{utterance}  

    & & & \multicolumn{2}{p{0.3\linewidth}}{\cellcolor[rgb]{0.95,0.95,0.95}{%
	\makecell[{{p{\linewidth}}}]{% 
	  \tt {\tiny [GM$|$GM]}  
	 Strafpunkt: (4,2) ist nicht leer, sondern enthält '│'.\\ \tt \\ \tt Ihr habt gemeinsam 13 von maximal 16 Strafpunkten angesammelt.\\ \tt Bitte versuche es noch mal! 
	  } 
	   } 
	   } 
	 & & \\ 
 

    \theutterance \stepcounter{utterance}  

    & & \multicolumn{4}{p{0.6\linewidth}}{\cellcolor[rgb]{0.9,0.9,0.9}{%
	\makecell[{{p{\linewidth}}}]{% 
	  \tt {\tiny [GM$\rangle$B]}  
	 Strafpunkt: (4,2) ist nicht leer, sondern enthält '│'.\\ \tt \\ \tt Ihr habt gemeinsam 13 von maximal 16 Strafpunkten angesammelt.\\ \tt Bitte versuche es noch mal! 
	  } 
	   } 
	   } 
	 & \\ 
 

    \theutterance \stepcounter{utterance}  

    & & & \multicolumn{4}{p{0.6\linewidth}}{\cellcolor[rgb]{1,0.85,0.72}{%
	\makecell[{{p{\linewidth}}}]{% 
	  \tt {\tiny [GM$\langle$B]}  
	 VERSCHIEBE: T, (3,3) 
	  } 
	   } 
	   } 
	  \\ 
 

    \theutterance \stepcounter{utterance}  

    & & & \multicolumn{2}{p{0.3\linewidth}}{\cellcolor[rgb]{0.95,0.95,0.95}{%
	\makecell[{{p{\linewidth}}}]{% 
	  \tt {\tiny [GM$|$GM]}  
	 VERSCHIEBE: T, (3,3) 
	  } 
	   } 
	   } 
	 & & \\ 
 

    \theutterance \stepcounter{utterance}  

    & & & \multicolumn{2}{p{0.3\linewidth}}{\cellcolor[rgb]{0.95,0.95,0.95}{%
	\makecell[{{p{\linewidth}}}]{% 
	  \tt {\tiny [GM$|$GM]}  
	 'T' erfolgreich nach (3,3) verschoben. Dein Raster sieht jetzt so aus:\\ \tt ```\\ \tt  1234567\\ \tt ╔═══╤═╤═╗\\ \tt ║PM◌│S│◌║ 1\\ \tt ║D◌O│◌├─╢ 2\\ \tt ║◌◌T├─┘◌║ 3\\ \tt ║◌◌◌│◌◌◌║ 4\\ \tt ╟─┬─┼───╢ 5\\ \tt ║◌│◌│◌◌◌║ 6\\ \tt ║◌│◌│◌◌A║ 7\\ \tt ╚═╧═╧═══╝\\ \tt ```\\ \tt  
	  } 
	   } 
	   } 
	 & & \\ 
 

    \theutterance \stepcounter{utterance}  

    & & \multicolumn{4}{p{0.6\linewidth}}{\cellcolor[rgb]{0.9,0.9,0.9}{%
	\makecell[{{p{\linewidth}}}]{% 
	  \tt {\tiny [A$\langle$GM]}  
	 'T' erfolgreich nach (5,2) verschoben. Dein Raster sieht jetzt so aus:\\ \tt ```\\ \tt  1234567\\ \tt ╔═══╤═╤═╗\\ \tt ║P◌◌│◌│◌║ 1\\ \tt ║A◌◌│T├─╢ 2\\ \tt ║◌◌M├─┘◌║ 3\\ \tt ║◌◌◌│◌◌◌║ 4\\ \tt ╟─┬─┼───╢ 5\\ \tt ║◌│◌│S◌O║ 6\\ \tt ║◌│◌│◌D◌║ 7\\ \tt ╚═╧═╧═══╝\\ \tt ```\\ \tt Ihr seid gerade bei Runde 11 von maximal 28.\\ \tt Ihr habt gemeinsam 13 von maximal 16 Strafpunkten angesammelt.\\ \tt Dein Mitspieler hat ein Objekt auf seinem Raster bewegt.\\ \tt \\ \tt \\ \tt Wie lautet dein nächster Befehl? 
	  } 
	   } 
	   } 
	 & \\ 
 

    \theutterance \stepcounter{utterance}  

    & \multicolumn{4}{p{0.6\linewidth}}{\cellcolor[rgb]{0.8,1,0.9}{%
	\makecell[{{p{\linewidth}}}]{% 
	  \tt {\tiny [A$\rangle$GM]}  
	 SPRICH: T ist jetzt auf (5,2), nicht (4,2). Bitte aktualisiere deine Zielposition für T auf (5,2). Ich bewege nun S. 
	  } 
	   } 
	   } 
	 & & \\ 
 

    \theutterance \stepcounter{utterance}  

    & & & \multicolumn{2}{p{0.3\linewidth}}{\cellcolor[rgb]{0.95,0.95,0.95}{%
	\makecell[{{p{\linewidth}}}]{% 
	  \tt {\tiny [GM$|$GM]}  
	 SPRICH: T ist jetzt auf (5,2), nicht (4,2). Bitte aktualisiere deine Zielposition für T auf (5,2). Ich bewege nun S. 
	  } 
	   } 
	   } 
	 & & \\ 
 

    \theutterance \stepcounter{utterance}  

    & & \multicolumn{4}{p{0.6\linewidth}}{\cellcolor[rgb]{0.9,0.9,0.9}{%
	\makecell[{{p{\linewidth}}}]{% 
	  \tt {\tiny [GM$\rangle$B]}  
	 'T' erfolgreich nach (3,3) verschoben. Dein Raster sieht jetzt so aus:\\ \tt ```\\ \tt  1234567\\ \tt ╔═══╤═╤═╗\\ \tt ║PM◌│S│◌║ 1\\ \tt ║D◌O│◌├─╢ 2\\ \tt ║◌◌T├─┘◌║ 3\\ \tt ║◌◌◌│◌◌◌║ 4\\ \tt ╟─┬─┼───╢ 5\\ \tt ║◌│◌│◌◌◌║ 6\\ \tt ║◌│◌│◌◌A║ 7\\ \tt ╚═╧═╧═══╝\\ \tt ```\\ \tt Ihr seid gerade bei Runde 12 von maximal 28.\\ \tt Ihr habt gemeinsam 13 von maximal 16 Strafpunkten angesammelt.\\ \tt Dein Mitspieler hat dir diese Nachricht geschickt:\\ \tt \\ \tt \\ \tt "T ist jetzt auf (5,2), nicht (4,2). Bitte aktualisiere deine Zielposition für T auf (5,2). Ich bewege nun S."\\ \tt \\ \tt \\ \tt Wie lautet dein nächster Befehl? 
	  } 
	   } 
	   } 
	 & \\ 
 

    \theutterance \stepcounter{utterance}  

    & & & \multicolumn{4}{p{0.6\linewidth}}{\cellcolor[rgb]{1,0.85,0.72}{%
	\makecell[{{p{\linewidth}}}]{% 
	  \tt {\tiny [GM$\langle$B]}  
	 SPRICH: Verstanden, T soll (5,2) sein. Ich bewege nun A nach (1,2) und dann M nach (3,3). Bitte bestätige. 
	  } 
	   } 
	   } 
	  \\ 
 

    \theutterance \stepcounter{utterance}  

    & & & \multicolumn{2}{p{0.3\linewidth}}{\cellcolor[rgb]{0.95,0.95,0.95}{%
	\makecell[{{p{\linewidth}}}]{% 
	  \tt {\tiny [GM$|$GM]}  
	 SPRICH: Verstanden, T soll (5,2) sein. Ich bewege nun A nach (1,2) und dann M nach (3,3). Bitte bestätige. 
	  } 
	   } 
	   } 
	 & & \\ 
 

    \theutterance \stepcounter{utterance}  

    & & \multicolumn{4}{p{0.6\linewidth}}{\cellcolor[rgb]{0.9,0.9,0.9}{%
	\makecell[{{p{\linewidth}}}]{% 
	  \tt {\tiny [A$\langle$GM]}  
	 Deine Nachricht wurde an deinen Mitspieler übermittelt.\\ \tt \\ \tt \\ \tt Ihr seid gerade bei Runde 12 von maximal 28.\\ \tt Ihr habt gemeinsam 13 von maximal 16 Strafpunkten angesammelt.\\ \tt Dein Mitspieler hat dir diese Nachricht geschickt:\\ \tt \\ \tt \\ \tt "Verstanden, T soll (5,2) sein. Ich bewege nun A nach (1,2) und dann M nach (3,3). Bitte bestätige."\\ \tt \\ \tt \\ \tt Wie lautet dein nächster Befehl? 
	  } 
	   } 
	   } 
	 & \\ 
 

    \theutterance \stepcounter{utterance}  

    & \multicolumn{4}{p{0.6\linewidth}}{\cellcolor[rgb]{0.8,1,0.9}{%
	\makecell[{{p{\linewidth}}}]{% 
	  \tt {\tiny [A$\rangle$GM]}  
	 SPRICH: A ist bereits auf (1,2) bei mir. M ist auf (3,3). Ich bewege nun S nach (4,1). Bitte bestätige. 
	  } 
	   } 
	   } 
	 & & \\ 
 

    \theutterance \stepcounter{utterance}  

    & & & \multicolumn{2}{p{0.3\linewidth}}{\cellcolor[rgb]{0.95,0.95,0.95}{%
	\makecell[{{p{\linewidth}}}]{% 
	  \tt {\tiny [GM$|$GM]}  
	 SPRICH: A ist bereits auf (1,2) bei mir. M ist auf (3,3). Ich bewege nun S nach (4,1). Bitte bestätige. 
	  } 
	   } 
	   } 
	 & & \\ 
 

    \theutterance \stepcounter{utterance}  

    & & \multicolumn{4}{p{0.6\linewidth}}{\cellcolor[rgb]{0.9,0.9,0.9}{%
	\makecell[{{p{\linewidth}}}]{% 
	  \tt {\tiny [GM$\rangle$B]}  
	 Deine Nachricht wurde an deinen Mitspieler übermittelt.\\ \tt \\ \tt \\ \tt Ihr seid gerade bei Runde 13 von maximal 28.\\ \tt Ihr habt gemeinsam 13 von maximal 16 Strafpunkten angesammelt.\\ \tt Dein Mitspieler hat dir diese Nachricht geschickt:\\ \tt \\ \tt \\ \tt "A ist bereits auf (1,2) bei mir. M ist auf (3,3). Ich bewege nun S nach (4,1). Bitte bestätige."\\ \tt \\ \tt \\ \tt Wie lautet dein nächster Befehl? 
	  } 
	   } 
	   } 
	 & \\ 
 

    \theutterance \stepcounter{utterance}  

    & & & \multicolumn{4}{p{0.6\linewidth}}{\cellcolor[rgb]{1,0.85,0.72}{%
	\makecell[{{p{\linewidth}}}]{% 
	  \tt {\tiny [GM$\langle$B]}  
	 SPRICH: S ist bereits auf (4,1) bei mir. A muss noch nach (1,2) bewegt werden. Ich bewege A jetzt. 
	  } 
	   } 
	   } 
	  \\ 
 

    \theutterance \stepcounter{utterance}  

    & & & \multicolumn{2}{p{0.3\linewidth}}{\cellcolor[rgb]{0.95,0.95,0.95}{%
	\makecell[{{p{\linewidth}}}]{% 
	  \tt {\tiny [GM$|$GM]}  
	 SPRICH: S ist bereits auf (4,1) bei mir. A muss noch nach (1,2) bewegt werden. Ich bewege A jetzt. 
	  } 
	   } 
	   } 
	 & & \\ 
 

    \theutterance \stepcounter{utterance}  

    & & \multicolumn{4}{p{0.6\linewidth}}{\cellcolor[rgb]{0.9,0.9,0.9}{%
	\makecell[{{p{\linewidth}}}]{% 
	  \tt {\tiny [A$\langle$GM]}  
	 Deine Nachricht wurde an deinen Mitspieler übermittelt.\\ \tt \\ \tt \\ \tt Ihr seid gerade bei Runde 13 von maximal 28.\\ \tt Ihr habt gemeinsam 13 von maximal 16 Strafpunkten angesammelt.\\ \tt Dein Mitspieler hat dir diese Nachricht geschickt:\\ \tt \\ \tt \\ \tt "S ist bereits auf (4,1) bei mir. A muss noch nach (1,2) bewegt werden. Ich bewege A jetzt."\\ \tt \\ \tt \\ \tt Wie lautet dein nächster Befehl? 
	  } 
	   } 
	   } 
	 & \\ 
 

    \theutterance \stepcounter{utterance}  

    & \multicolumn{4}{p{0.6\linewidth}}{\cellcolor[rgb]{0.8,1,0.9}{%
	\makecell[{{p{\linewidth}}}]{% 
	  \tt {\tiny [A$\rangle$GM]}  
	 VERSCHIEBE: S, (4, 1) 
	  } 
	   } 
	   } 
	 & & \\ 
 

    \theutterance \stepcounter{utterance}  

    & & & \multicolumn{2}{p{0.3\linewidth}}{\cellcolor[rgb]{0.95,0.95,0.95}{%
	\makecell[{{p{\linewidth}}}]{% 
	  \tt {\tiny [GM$|$GM]}  
	 VERSCHIEBE: S, (4, 1) 
	  } 
	   } 
	   } 
	 & & \\ 
 

    \theutterance \stepcounter{utterance}  

    & & & \multicolumn{2}{p{0.3\linewidth}}{\cellcolor[rgb]{0.95,0.95,0.95}{%
	\makecell[{{p{\linewidth}}}]{% 
	  \tt {\tiny [GM$|$GM]}  
	 Strafpunkt: (4,1) ist nicht leer, sondern enthält '│'.\\ \tt \\ \tt Ihr habt gemeinsam 14 von maximal 16 Strafpunkten angesammelt.\\ \tt Bitte versuche es noch mal! 
	  } 
	   } 
	   } 
	 & & \\ 
 

    \theutterance \stepcounter{utterance}  

    & & \multicolumn{4}{p{0.6\linewidth}}{\cellcolor[rgb]{0.9,0.9,0.9}{%
	\makecell[{{p{\linewidth}}}]{% 
	  \tt {\tiny [A$\langle$GM]}  
	 Strafpunkt: (4,1) ist nicht leer, sondern enthält '│'.\\ \tt \\ \tt Ihr habt gemeinsam 14 von maximal 16 Strafpunkten angesammelt.\\ \tt Bitte versuche es noch mal! 
	  } 
	   } 
	   } 
	 & \\ 
 

    \theutterance \stepcounter{utterance}  

    & \multicolumn{4}{p{0.6\linewidth}}{\cellcolor[rgb]{0.8,1,0.9}{%
	\makecell[{{p{\linewidth}}}]{% 
	  \tt {\tiny [A$\rangle$GM]}  
	 SPRICH: (4,1) ist bei mir eine Trennlinie. Sollte S nach (4,1)? Bitte alternative Position für S vorschlagen. 
	  } 
	   } 
	   } 
	 & & \\ 
 

    \theutterance \stepcounter{utterance}  

    & & & \multicolumn{2}{p{0.3\linewidth}}{\cellcolor[rgb]{0.95,0.95,0.95}{%
	\makecell[{{p{\linewidth}}}]{% 
	  \tt {\tiny [GM$|$GM]}  
	 SPRICH: (4,1) ist bei mir eine Trennlinie. Sollte S nach (4,1)? Bitte alternative Position für S vorschlagen. 
	  } 
	   } 
	   } 
	 & & \\ 
 

    \theutterance \stepcounter{utterance}  

    & & \multicolumn{4}{p{0.6\linewidth}}{\cellcolor[rgb]{0.9,0.9,0.9}{%
	\makecell[{{p{\linewidth}}}]{% 
	  \tt {\tiny [GM$\rangle$B]}  
	 Deine Nachricht wurde an deinen Mitspieler übermittelt.\\ \tt \\ \tt \\ \tt Ihr seid gerade bei Runde 14 von maximal 28.\\ \tt Ihr habt gemeinsam 14 von maximal 16 Strafpunkten angesammelt.\\ \tt Dein Mitspieler hat dir diese Nachricht geschickt:\\ \tt \\ \tt \\ \tt "(4,1) ist bei mir eine Trennlinie. Sollte S nach (4,1)? Bitte alternative Position für S vorschlagen."\\ \tt \\ \tt \\ \tt Wie lautet dein nächster Befehl? 
	  } 
	   } 
	   } 
	 & \\ 
 

    \theutterance \stepcounter{utterance}  

    & & & \multicolumn{4}{p{0.6\linewidth}}{\cellcolor[rgb]{1,0.85,0.72}{%
	\makecell[{{p{\linewidth}}}]{% 
	  \tt {\tiny [GM$\langle$B]}  
	 SPRICH: Auf meinem Brett ist (4,1) leer (◌). Wenn es bei dir eine Trennlinie ist, schlage (5,1) für S vor. Bitte prüfen. 
	  } 
	   } 
	   } 
	  \\ 
 

    \theutterance \stepcounter{utterance}  

    & & & \multicolumn{2}{p{0.3\linewidth}}{\cellcolor[rgb]{0.95,0.95,0.95}{%
	\makecell[{{p{\linewidth}}}]{% 
	  \tt {\tiny [GM$|$GM]}  
	 SPRICH: Auf meinem Brett ist (4,1) leer (◌). Wenn es bei dir eine Trennlinie ist, schlage (5,1) für S vor. Bitte prüfen. 
	  } 
	   } 
	   } 
	 & & \\ 
 

    \theutterance \stepcounter{utterance}  

    & & \multicolumn{4}{p{0.6\linewidth}}{\cellcolor[rgb]{0.9,0.9,0.9}{%
	\makecell[{{p{\linewidth}}}]{% 
	  \tt {\tiny [A$\langle$GM]}  
	 Deine Nachricht wurde an deinen Mitspieler übermittelt.\\ \tt \\ \tt \\ \tt Ihr seid gerade bei Runde 14 von maximal 28.\\ \tt Ihr habt gemeinsam 14 von maximal 16 Strafpunkten angesammelt.\\ \tt Dein Mitspieler hat dir diese Nachricht geschickt:\\ \tt \\ \tt \\ \tt "Auf meinem Brett ist (4,1) leer (◌). Wenn es bei dir eine Trennlinie ist, schlage (5,1) für S vor. Bitte prüfen."\\ \tt \\ \tt \\ \tt Wie lautet dein nächster Befehl? 
	  } 
	   } 
	   } 
	 & \\ 
 

    \theutterance \stepcounter{utterance}  

    & \multicolumn{4}{p{0.6\linewidth}}{\cellcolor[rgb]{0.8,1,0.9}{%
	\makecell[{{p{\linewidth}}}]{% 
	  \tt {\tiny [A$\rangle$GM]}  
	 VERSCHIEBE: S, (5, 1) 
	  } 
	   } 
	   } 
	 & & \\ 
 

    \theutterance \stepcounter{utterance}  

    & & & \multicolumn{2}{p{0.3\linewidth}}{\cellcolor[rgb]{0.95,0.95,0.95}{%
	\makecell[{{p{\linewidth}}}]{% 
	  \tt {\tiny [GM$|$GM]}  
	 VERSCHIEBE: S, (5, 1) 
	  } 
	   } 
	   } 
	 & & \\ 
 

    \theutterance \stepcounter{utterance}  

    & & & \multicolumn{2}{p{0.3\linewidth}}{\cellcolor[rgb]{0.95,0.95,0.95}{%
	\makecell[{{p{\linewidth}}}]{% 
	  \tt {\tiny [GM$|$GM]}  
	 'S' erfolgreich nach (5,1) verschoben. Dein Raster sieht jetzt so aus:\\ \tt ```\\ \tt  1234567\\ \tt ╔═══╤═╤═╗\\ \tt ║P◌◌│S│◌║ 1\\ \tt ║A◌◌│T├─╢ 2\\ \tt ║◌◌M├─┘◌║ 3\\ \tt ║◌◌◌│◌◌◌║ 4\\ \tt ╟─┬─┼───╢ 5\\ \tt ║◌│◌│◌◌O║ 6\\ \tt ║◌│◌│◌D◌║ 7\\ \tt ╚═╧═╧═══╝\\ \tt ```\\ \tt  
	  } 
	   } 
	   } 
	 & & \\ 
 

    \theutterance \stepcounter{utterance}  

    & & \multicolumn{4}{p{0.6\linewidth}}{\cellcolor[rgb]{0.9,0.9,0.9}{%
	\makecell[{{p{\linewidth}}}]{% 
	  \tt {\tiny [GM$\rangle$B]}  
	 Deine Nachricht wurde an deinen Mitspieler übermittelt.\\ \tt \\ \tt \\ \tt Ihr seid gerade bei Runde 15 von maximal 28.\\ \tt Ihr habt gemeinsam 14 von maximal 16 Strafpunkten angesammelt.\\ \tt Dein Mitspieler hat ein Objekt auf seinem Raster bewegt.\\ \tt \\ \tt \\ \tt Wie lautet dein nächster Befehl? 
	  } 
	   } 
	   } 
	 & \\ 
 

    \theutterance \stepcounter{utterance}  

    & & & \multicolumn{4}{p{0.6\linewidth}}{\cellcolor[rgb]{1,0.85,0.72}{%
	\makecell[{{p{\linewidth}}}]{% 
	  \tt {\tiny [GM$\langle$B]}  
	 VERSCHIEBE: A, (1,2) 
	  } 
	   } 
	   } 
	  \\ 
 

    \theutterance \stepcounter{utterance}  

    & & & \multicolumn{2}{p{0.3\linewidth}}{\cellcolor[rgb]{0.95,0.95,0.95}{%
	\makecell[{{p{\linewidth}}}]{% 
	  \tt {\tiny [GM$|$GM]}  
	 VERSCHIEBE: A, (1,2) 
	  } 
	   } 
	   } 
	 & & \\ 
 

    \theutterance \stepcounter{utterance}  

    & & & \multicolumn{2}{p{0.3\linewidth}}{\cellcolor[rgb]{0.95,0.95,0.95}{%
	\makecell[{{p{\linewidth}}}]{% 
	  \tt {\tiny [GM$|$GM]}  
	 Strafpunkt: (1,2) ist nicht leer, sondern enthält 'D'.\\ \tt \\ \tt Ihr habt gemeinsam 15 von maximal 16 Strafpunkten angesammelt.\\ \tt Bitte versuche es noch mal! 
	  } 
	   } 
	   } 
	 & & \\ 
 

    \theutterance \stepcounter{utterance}  

    & & \multicolumn{4}{p{0.6\linewidth}}{\cellcolor[rgb]{0.9,0.9,0.9}{%
	\makecell[{{p{\linewidth}}}]{% 
	  \tt {\tiny [GM$\rangle$B]}  
	 Strafpunkt: (1,2) ist nicht leer, sondern enthält 'D'.\\ \tt \\ \tt Ihr habt gemeinsam 15 von maximal 16 Strafpunkten angesammelt.\\ \tt Bitte versuche es noch mal! 
	  } 
	   } 
	   } 
	 & \\ 
 

    \theutterance \stepcounter{utterance}  

    & & & \multicolumn{4}{p{0.6\linewidth}}{\cellcolor[rgb]{1,0.85,0.72}{%
	\makecell[{{p{\linewidth}}}]{% 
	  \tt {\tiny [GM$\langle$B]}  
	 VERSCHIEBE: A, (2,2) 
	  } 
	   } 
	   } 
	  \\ 
 

    \theutterance \stepcounter{utterance}  

    & & & \multicolumn{2}{p{0.3\linewidth}}{\cellcolor[rgb]{0.95,0.95,0.95}{%
	\makecell[{{p{\linewidth}}}]{% 
	  \tt {\tiny [GM$|$GM]}  
	 VERSCHIEBE: A, (2,2) 
	  } 
	   } 
	   } 
	 & & \\ 
 

    \theutterance \stepcounter{utterance}  

    & & & \multicolumn{2}{p{0.3\linewidth}}{\cellcolor[rgb]{0.95,0.95,0.95}{%
	\makecell[{{p{\linewidth}}}]{% 
	  \tt {\tiny [GM$|$GM]}  
	 'A' erfolgreich nach (2,2) verschoben. Dein Raster sieht jetzt so aus:\\ \tt ```\\ \tt  1234567\\ \tt ╔═══╤═╤═╗\\ \tt ║PM◌│S│◌║ 1\\ \tt ║DAO│◌├─╢ 2\\ \tt ║◌◌T├─┘◌║ 3\\ \tt ║◌◌◌│◌◌◌║ 4\\ \tt ╟─┬─┼───╢ 5\\ \tt ║◌│◌│◌◌◌║ 6\\ \tt ║◌│◌│◌◌◌║ 7\\ \tt ╚═╧═╧═══╝\\ \tt ```\\ \tt  
	  } 
	   } 
	   } 
	 & & \\ 
 

    \theutterance \stepcounter{utterance}  

    & & \multicolumn{4}{p{0.6\linewidth}}{\cellcolor[rgb]{0.9,0.9,0.9}{%
	\makecell[{{p{\linewidth}}}]{% 
	  \tt {\tiny [A$\langle$GM]}  
	 'S' erfolgreich nach (5,1) verschoben. Dein Raster sieht jetzt so aus:\\ \tt ```\\ \tt  1234567\\ \tt ╔═══╤═╤═╗\\ \tt ║P◌◌│S│◌║ 1\\ \tt ║A◌◌│T├─╢ 2\\ \tt ║◌◌M├─┘◌║ 3\\ \tt ║◌◌◌│◌◌◌║ 4\\ \tt ╟─┬─┼───╢ 5\\ \tt ║◌│◌│◌◌O║ 6\\ \tt ║◌│◌│◌D◌║ 7\\ \tt ╚═╧═╧═══╝\\ \tt ```\\ \tt Ihr seid gerade bei Runde 15 von maximal 28.\\ \tt Ihr habt gemeinsam 15 von maximal 16 Strafpunkten angesammelt.\\ \tt Dein Mitspieler hat ein Objekt auf seinem Raster bewegt.\\ \tt \\ \tt \\ \tt Wie lautet dein nächster Befehl? 
	  } 
	   } 
	   } 
	 & \\ 
 

    \theutterance \stepcounter{utterance}  

    & \multicolumn{4}{p{0.6\linewidth}}{\cellcolor[rgb]{0.8,1,0.9}{%
	\makecell[{{p{\linewidth}}}]{% 
	  \tt {\tiny [A$\rangle$GM]}  
	 SPRICH: S ist jetzt auf (5,1) bei mir. Bitte aktualisiere Ziel für S auf (5,1). Ich bewege nun D nach (5,1) - ist das frei? 
	  } 
	   } 
	   } 
	 & & \\ 
 

    \theutterance \stepcounter{utterance}  

    & & & \multicolumn{2}{p{0.3\linewidth}}{\cellcolor[rgb]{0.95,0.95,0.95}{%
	\makecell[{{p{\linewidth}}}]{% 
	  \tt {\tiny [GM$|$GM]}  
	 SPRICH: S ist jetzt auf (5,1) bei mir. Bitte aktualisiere Ziel für S auf (5,1). Ich bewege nun D nach (5,1) - ist das frei? 
	  } 
	   } 
	   } 
	 & & \\ 
 

    \theutterance \stepcounter{utterance}  

    & & \multicolumn{4}{p{0.6\linewidth}}{\cellcolor[rgb]{0.9,0.9,0.9}{%
	\makecell[{{p{\linewidth}}}]{% 
	  \tt {\tiny [GM$\rangle$B]}  
	 'A' erfolgreich nach (2,2) verschoben. Dein Raster sieht jetzt so aus:\\ \tt ```\\ \tt  1234567\\ \tt ╔═══╤═╤═╗\\ \tt ║PM◌│S│◌║ 1\\ \tt ║DAO│◌├─╢ 2\\ \tt ║◌◌T├─┘◌║ 3\\ \tt ║◌◌◌│◌◌◌║ 4\\ \tt ╟─┬─┼───╢ 5\\ \tt ║◌│◌│◌◌◌║ 6\\ \tt ║◌│◌│◌◌◌║ 7\\ \tt ╚═╧═╧═══╝\\ \tt ```\\ \tt Ihr seid gerade bei Runde 16 von maximal 28.\\ \tt Ihr habt gemeinsam 15 von maximal 16 Strafpunkten angesammelt.\\ \tt Dein Mitspieler hat dir diese Nachricht geschickt:\\ \tt \\ \tt \\ \tt "S ist jetzt auf (5,1) bei mir. Bitte aktualisiere Ziel für S auf (5,1). Ich bewege nun D nach (5,1) - ist das frei?"\\ \tt \\ \tt \\ \tt Wie lautet dein nächster Befehl? 
	  } 
	   } 
	   } 
	 & \\ 
 

    \theutterance \stepcounter{utterance}  

    & & & \multicolumn{4}{p{0.6\linewidth}}{\cellcolor[rgb]{1,0.85,0.72}{%
	\makecell[{{p{\linewidth}}}]{% 
	  \tt {\tiny [GM$\langle$B]}  
	 SPRICH: (5,1) ist bei mir belegt mit S. Bitte D nach (1,1) oder (3,4) bewegen. Prüfe diese. 
	  } 
	   } 
	   } 
	  \\ 
 

    \theutterance \stepcounter{utterance}  

    & & & \multicolumn{2}{p{0.3\linewidth}}{\cellcolor[rgb]{0.95,0.95,0.95}{%
	\makecell[{{p{\linewidth}}}]{% 
	  \tt {\tiny [GM$|$GM]}  
	 SPRICH: (5,1) ist bei mir belegt mit S. Bitte D nach (1,1) oder (3,4) bewegen. Prüfe diese. 
	  } 
	   } 
	   } 
	 & & \\ 
 

    \theutterance \stepcounter{utterance}  

    & & \multicolumn{4}{p{0.6\linewidth}}{\cellcolor[rgb]{0.9,0.9,0.9}{%
	\makecell[{{p{\linewidth}}}]{% 
	  \tt {\tiny [A$\langle$GM]}  
	 Deine Nachricht wurde an deinen Mitspieler übermittelt.\\ \tt \\ \tt \\ \tt Ihr seid gerade bei Runde 16 von maximal 28.\\ \tt Ihr habt gemeinsam 15 von maximal 16 Strafpunkten angesammelt.\\ \tt Dein Mitspieler hat dir diese Nachricht geschickt:\\ \tt \\ \tt \\ \tt "(5,1) ist bei mir belegt mit S. Bitte D nach (1,1) oder (3,4) bewegen. Prüfe diese."\\ \tt \\ \tt \\ \tt Wie lautet dein nächster Befehl? 
	  } 
	   } 
	   } 
	 & \\ 
 

    \theutterance \stepcounter{utterance}  

    & \multicolumn{4}{p{0.6\linewidth}}{\cellcolor[rgb]{0.8,1,0.9}{%
	\makecell[{{p{\linewidth}}}]{% 
	  \tt {\tiny [A$\rangle$GM]}  
	 VERSCHIEBE: D, (1, 1) 
	  } 
	   } 
	   } 
	 & & \\ 
 

    \theutterance \stepcounter{utterance}  

    & & & \multicolumn{2}{p{0.3\linewidth}}{\cellcolor[rgb]{0.95,0.95,0.95}{%
	\makecell[{{p{\linewidth}}}]{% 
	  \tt {\tiny [GM$|$GM]}  
	 VERSCHIEBE: D, (1, 1) 
	  } 
	   } 
	   } 
	 & & \\ 
 

    \theutterance \stepcounter{utterance}  

    & & & \multicolumn{2}{p{0.3\linewidth}}{\cellcolor[rgb]{0.95,0.95,0.95}{%
	\makecell[{{p{\linewidth}}}]{% 
	  \tt {\tiny [GM$|$GM]}  
	 Strafpunkt: (1,1) ist nicht leer, sondern enthält 'P'.\\ \tt \\ \tt Ihr habt gemeinsam 16 von maximal 16 Strafpunkten angesammelt.\\ \tt Bitte versuche es noch mal! 
	  } 
	   } 
	   } 
	 & & \\ 
 

    \theutterance \stepcounter{utterance}  

    & & \multicolumn{4}{p{0.6\linewidth}}{\cellcolor[rgb]{0.9,0.9,0.9}{%
	\makecell[{{p{\linewidth}}}]{% 
	  \tt {\tiny [A$\langle$GM]}  
	 Strafpunkt: (1,1) ist nicht leer, sondern enthält 'P'.\\ \tt \\ \tt Ihr habt gemeinsam 16 von maximal 16 Strafpunkten angesammelt.\\ \tt Bitte versuche es noch mal! 
	  } 
	   } 
	   } 
	 & \\ 
 

    \theutterance \stepcounter{utterance}  

    & \multicolumn{4}{p{0.6\linewidth}}{\cellcolor[rgb]{0.8,1,0.9}{%
	\makecell[{{p{\linewidth}}}]{% 
	  \tt {\tiny [A$\rangle$GM]}  
	 SPRICH: Wir haben verloren wegen Strafpunkten. Können wir neu starten oder ist das Spiel vorbei? 
	  } 
	   } 
	   } 
	 & & \\ 
 

    \theutterance \stepcounter{utterance}  

    & & & \multicolumn{2}{p{0.3\linewidth}}{\cellcolor[rgb]{0.95,0.95,0.95}{%
	\makecell[{{p{\linewidth}}}]{% 
	  \tt {\tiny [GM$|$GM]}  
	 SPRICH: Wir haben verloren wegen Strafpunkten. Können wir neu starten oder ist das Spiel vorbei? 
	  } 
	   } 
	   } 
	 & & \\ 
 

    \theutterance \stepcounter{utterance}  

    & & \multicolumn{4}{p{0.6\linewidth}}{\cellcolor[rgb]{0.9,0.9,0.9}{%
	\makecell[{{p{\linewidth}}}]{% 
	  \tt {\tiny [GM$\rangle$B]}  
	 Deine Nachricht wurde an deinen Mitspieler übermittelt.\\ \tt \\ \tt \\ \tt Ihr seid gerade bei Runde 17 von maximal 28.\\ \tt Ihr habt gemeinsam 16 von maximal 16 Strafpunkten angesammelt.\\ \tt Dein Mitspieler hat dir diese Nachricht geschickt:\\ \tt \\ \tt \\ \tt "Wir haben verloren wegen Strafpunkten. Können wir neu starten oder ist das Spiel vorbei?"\\ \tt \\ \tt \\ \tt Wie lautet dein nächster Befehl? 
	  } 
	   } 
	   } 
	 & \\ 
 

    \theutterance \stepcounter{utterance}  

    & & & \multicolumn{4}{p{0.6\linewidth}}{\cellcolor[rgb]{1,0.85,0.72}{%
	\makecell[{{p{\linewidth}}}]{% 
	  \tt {\tiny [GM$\langle$B]}  
	 SPRICH: Das Spiel ist verloren wegen 16 Strafpunkten. Wir sollten neu starten und besser koordinieren, z.B. zuerst alle freien Positionen abgleichen. 
	  } 
	   } 
	   } 
	  \\ 
 

    \theutterance \stepcounter{utterance}  

    & & & \multicolumn{2}{p{0.3\linewidth}}{\cellcolor[rgb]{0.95,0.95,0.95}{%
	\makecell[{{p{\linewidth}}}]{% 
	  \tt {\tiny [GM$|$GM]}  
	 SPRICH: Das Spiel ist verloren wegen 16 Strafpunkten. Wir sollten neu starten und besser koordinieren, z.B. zuerst alle freien Positionen abgleichen. 
	  } 
	   } 
	   } 
	 & & \\ 
 

    \theutterance \stepcounter{utterance}  

    & & \multicolumn{4}{p{0.6\linewidth}}{\cellcolor[rgb]{0.9,0.9,0.9}{%
	\makecell[{{p{\linewidth}}}]{% 
	  \tt {\tiny [A$\langle$GM]}  
	 Deine Nachricht wurde an deinen Mitspieler übermittelt.\\ \tt \\ \tt \\ \tt Ihr seid gerade bei Runde 17 von maximal 28.\\ \tt Ihr habt gemeinsam 16 von maximal 16 Strafpunkten angesammelt.\\ \tt Dein Mitspieler hat dir diese Nachricht geschickt:\\ \tt \\ \tt \\ \tt "Das Spiel ist verloren wegen 16 Strafpunkten. Wir sollten neu starten und besser koordinieren, z.B. zuerst alle freien Positionen abgleichen."\\ \tt \\ \tt \\ \tt Wie lautet dein nächster Befehl? 
	  } 
	   } 
	   } 
	 & \\ 
 

    \theutterance \stepcounter{utterance}  

    & \multicolumn{4}{p{0.6\linewidth}}{\cellcolor[rgb]{0.8,1,0.9}{%
	\makecell[{{p{\linewidth}}}]{% 
	  \tt {\tiny [A$\rangle$GM]}  
	 SPRICH: Ja, lassen wir neu starten. Ich beginne mit meinen freien Positionen: (1,4), (2,4), (3,4), (4,4), (5,4), (6,4), (7,4), (1,6), (2,6), (3,6), (4,6), (5,6), (6,6), (7,6), (1,7), (2,7), (3,7), (4,7), (5,7), (6,7), (7,7). Welche sind bei dir frei? 
	  } 
	   } 
	   } 
	 & & \\ 
 

    \theutterance \stepcounter{utterance}  

    & & & \multicolumn{2}{p{0.3\linewidth}}{\cellcolor[rgb]{0.95,0.95,0.95}{%
	\makecell[{{p{\linewidth}}}]{% 
	  \tt {\tiny [GM$|$GM]}  
	 SPRICH: Ja, lassen wir neu starten. Ich beginne mit meinen freien Positionen: (1,4), (2,4), (3,4), (4,4), (5,4), (6,4), (7,4), (1,6), (2,6), (3,6), (4,6), (5,6), (6,6), (7,6), (1,7), (2,7), (3,7), (4,7), (5,7), (6,7), (7,7). Welche sind bei dir frei? 
	  } 
	   } 
	   } 
	 & & \\ 
 

    \theutterance \stepcounter{utterance}  

    & & \multicolumn{4}{p{0.6\linewidth}}{\cellcolor[rgb]{0.9,0.9,0.9}{%
	\makecell[{{p{\linewidth}}}]{% 
	  \tt {\tiny [GM$\rangle$B]}  
	 Deine Nachricht wurde an deinen Mitspieler übermittelt.\\ \tt \\ \tt \\ \tt Ihr seid gerade bei Runde 18 von maximal 28.\\ \tt Ihr habt gemeinsam 16 von maximal 16 Strafpunkten angesammelt.\\ \tt Dein Mitspieler hat dir diese Nachricht geschickt:\\ \tt \\ \tt \\ \tt "Ja, lassen wir neu starten. Ich beginne mit meinen freien Positionen: (1,4), (2,4), (3,4), (4,4), (5,4), (6,4), (7,4), (1,6), (2,6), (3,6), (4,6), (5,6), (6,6), (7,6), (1,7), (2,7), (3,7), (4,7), (5,7), (6,7), (7,7). Welche sind bei dir frei?"\\ \tt \\ \tt \\ \tt Wie lautet dein nächster Befehl? 
	  } 
	   } 
	   } 
	 & \\ 
 

    \theutterance \stepcounter{utterance}  

    & & & \multicolumn{4}{p{0.6\linewidth}}{\cellcolor[rgb]{1,0.85,0.72}{%
	\makecell[{{p{\linewidth}}}]{% 
	  \tt {\tiny [GM$\langle$B]}  
	 SPRICH: Mein Raster hat nur 5 Spalten (1-5) und 7 Zeilen. Freie Positionen: (1,4), (2,4), (3,4), (4,4), (5,4), (1,5), (2,5), (3,5), (4,5), (5,5), (1,6), (2,6), (3,6), (4,6), (5,6), (1,7), (2,7), (3,7), (4,7), (5,7). Bitte wähle daraus. 
	  } 
	   } 
	   } 
	  \\ 
 

    \theutterance \stepcounter{utterance}  

    & & & \multicolumn{2}{p{0.3\linewidth}}{\cellcolor[rgb]{0.95,0.95,0.95}{%
	\makecell[{{p{\linewidth}}}]{% 
	  \tt {\tiny [GM$|$GM]}  
	 SPRICH: Mein Raster hat nur 5 Spalten (1-5) und 7 Zeilen. Freie Positionen: (1,4), (2,4), (3,4), (4,4), (5,4), (1,5), (2,5), (3,5), (4,5), (5,5), (1,6), (2,6), (3,6), (4,6), (5,6), (1,7), (2,7), (3,7), (4,7), (5,7). Bitte wähle daraus. 
	  } 
	   } 
	   } 
	 & & \\ 
 

    \theutterance \stepcounter{utterance}  

    & & \multicolumn{4}{p{0.6\linewidth}}{\cellcolor[rgb]{0.9,0.9,0.9}{%
	\makecell[{{p{\linewidth}}}]{% 
	  \tt {\tiny [A$\langle$GM]}  
	 Deine Nachricht wurde an deinen Mitspieler übermittelt.\\ \tt \\ \tt \\ \tt Ihr seid gerade bei Runde 18 von maximal 28.\\ \tt Ihr habt gemeinsam 16 von maximal 16 Strafpunkten angesammelt.\\ \tt Dein Mitspieler hat dir diese Nachricht geschickt:\\ \tt \\ \tt \\ \tt "Mein Raster hat nur 5 Spalten (1-5) und 7 Zeilen. Freie Positionen: (1,4), (2,4), (3,4), (4,4), (5,4), (1,5), (2,5), (3,5), (4,5), (5,5), (1,6), (2,6), (3,6), (4,6), (5,6), (1,7), (2,7), (3,7), (4,7), (5,7). Bitte wähle daraus."\\ \tt \\ \tt \\ \tt Wie lautet dein nächster Befehl? 
	  } 
	   } 
	   } 
	 & \\ 
 

    \theutterance \stepcounter{utterance}  

    & \multicolumn{4}{p{0.6\linewidth}}{\cellcolor[rgb]{0.8,1,0.9}{%
	\makecell[{{p{\linewidth}}}]{% 
	  \tt {\tiny [A$\rangle$GM]}  
	 SPRICH: Perfekt! Zielanordnung: P(1,4), O(2,4), T(3,4), S(4,4), D(5,4), A(1,5), M(2,5). Bitte bestätigen Sie alle Positionen. 
	  } 
	   } 
	   } 
	 & & \\ 
 

    \theutterance \stepcounter{utterance}  

    & & & \multicolumn{2}{p{0.3\linewidth}}{\cellcolor[rgb]{0.95,0.95,0.95}{%
	\makecell[{{p{\linewidth}}}]{% 
	  \tt {\tiny [GM$|$GM]}  
	 SPRICH: Perfekt! Zielanordnung: P(1,4), O(2,4), T(3,4), S(4,4), D(5,4), A(1,5), M(2,5). Bitte bestätigen Sie alle Positionen. 
	  } 
	   } 
	   } 
	 & & \\ 
 

    \theutterance \stepcounter{utterance}  

    & & \multicolumn{4}{p{0.6\linewidth}}{\cellcolor[rgb]{0.9,0.9,0.9}{%
	\makecell[{{p{\linewidth}}}]{% 
	  \tt {\tiny [GM$\rangle$B]}  
	 Deine Nachricht wurde an deinen Mitspieler übermittelt.\\ \tt \\ \tt \\ \tt Ihr seid gerade bei Runde 19 von maximal 28.\\ \tt Ihr habt gemeinsam 16 von maximal 16 Strafpunkten angesammelt.\\ \tt Dein Mitspieler hat dir diese Nachricht geschickt:\\ \tt \\ \tt \\ \tt "Perfekt! Zielanordnung: P(1,4), O(2,4), T(3,4), S(4,4), D(5,4), A(1,5), M(2,5). Bitte bestätigen Sie alle Positionen."\\ \tt \\ \tt \\ \tt Wie lautet dein nächster Befehl? 
	  } 
	   } 
	   } 
	 & \\ 
 

    \theutterance \stepcounter{utterance}  

    & & & \multicolumn{4}{p{0.6\linewidth}}{\cellcolor[rgb]{1,0.85,0.72}{%
	\makecell[{{p{\linewidth}}}]{% 
	  \tt {\tiny [GM$\langle$B]}  
	 Prüfe die vorgeschlagenen Positionen auf meinem aktuellen Raster:\\ \tt \\ \tt Aktuelles Raster:\\ \tt ```\\ \tt  1234567\\ \tt ╔═══╤═╤═╗\\ \tt ║PM◌│S│◌║ 1   // (1,1): P, (1,2): M, (1,3): ◌, (1,4): S, (1,5): ◌\\ \tt ║DAO│◌├─╢ 2   // (2,1): D, (2,2): A, (2,3): O, (2,4): ◌, (2,5): ─ (Trennlinie)\\ \tt ║◌◌T├─┘◌║ 3   // (3,1): ◌, (3,2): ◌, (3,3): T, (3,4): ─, (3,5): ◌\\ \tt ║◌◌◌│◌◌◌║ 4   // (4,1): ◌, (4,2): ◌, (4,3): ◌, (4,4): ◌, (4,5): ◌\\ \tt ╟─┬─┼───╢ 5   // Trennlinie\\ \tt ║◌│◌│◌◌◌║ 6   // (6,1): ◌, (6,2): ◌, (6,3): ◌, (6,4): ◌, (6,5): ◌\\ \tt ║◌│◌│◌◌◌║ 7   // (7,1): ◌, (7,2): ◌, (7,3): ◌, (7,4): ◌, (7,5): ◌\\ \tt ╚═╧═╧═══╝\\ \tt ```\\ \tt \\ \tt Vorgeschlagene Zielpositionen:\\ \tt - P(1,4): aktuell belegt mit S\\ \tt - O(2,4): aktuell leer (◌)\\ \tt - T(3,4): aktuell Trennlinie (─), ungültig\\ \tt - S(4,4): aktuell leer (◌)\\ \tt - D(5,4): ungültig (nur bis Spalte 5, aber (5,4) ist Trennlinie? Zeile 4 hat Spalte 4: ◌)\\ \tt - A(1,5): aktuell leer (◌)\\ \tt - M(2,5): aktuell Trennlinie (─), ungültig\\ \tt \\ \tt Korrektur: (5,4) existiert nicht, da Zeile 4 nur Spalten 1-5 hat. (5,4) wäre Zeile 4, Spalte 4? Aber Spalte 4 ist in Zeile 4 leer (◌). \\ \tt Tatsächlich:\\ \tt - (1,4): belegt (S)\\ \tt - (2,4): leer (◌)\\ \tt - (3,4): Trennlinie (─), ungültig\\ \tt - (4,4): leer (◌)\\ \tt - (5,4): nicht existent (nur Spalten 1-5)\\ \tt - (1,5): leer (◌)\\ \tt - (2,5): Trennlinie (─), ungültig\\ \tt \\ \tt Daher sind (3,4), (5,4), (2,5) ungültig. Ich muss dies mitteilen.\\ \tt \\ \tt SPRICH: (1,4) ist belegt, (3,4) und (2,5) sind Trennlinien, (5,4) existiert nicht. Bitte alternative Positionen vorschlagen, z.B. (4,4) für P, (1,5) für O, (3,5) für T, (4,5) für S, (5,5) für D, (1,6) für A, (2,6) für M. 
	  } 
	   } 
	   } 
	  \\ 
 

    \theutterance \stepcounter{utterance}  

    & & & \multicolumn{2}{p{0.3\linewidth}}{\cellcolor[rgb]{0.95,0.95,0.95}{%
	\makecell[{{p{\linewidth}}}]{% 
	  \tt {\tiny [GM$|$GM]}  
	 Prüfe die vorgeschlagenen Positionen auf meinem aktuellen Raster:\\ \tt \\ \tt Aktuelles Raster:\\ \tt ```\\ \tt  1234567\\ \tt ╔═══╤═╤═╗\\ \tt ║PM◌│S│◌║ 1   // (1,1): P, (1,2): M, (1,3): ◌, (1,4): S, (1,5): ◌\\ \tt ║DAO│◌├─╢ 2   // (2,1): D, (2,2): A, (2,3): O, (2,4): ◌, (2,5): ─ (Trennlinie)\\ \tt ║◌◌T├─┘◌║ 3   // (3,1): ◌, (3,2): ◌, (3,3): T, (3,4): ─, (3,5): ◌\\ \tt ║◌◌◌│◌◌◌║ 4   // (4,1): ◌, (4,2): ◌, (4,3): ◌, (4,4): ◌, (4,5): ◌\\ \tt ╟─┬─┼───╢ 5   // Trennlinie\\ \tt ║◌│◌│◌◌◌║ 6   // (6,1): ◌, (6,2): ◌, (6,3): ◌, (6,4): ◌, (6,5): ◌\\ \tt ║◌│◌│◌◌◌║ 7   // (7,1): ◌, (7,2): ◌, (7,3): ◌, (7,4): ◌, (7,5): ◌\\ \tt ╚═╧═╧═══╝\\ \tt ```\\ \tt \\ \tt Vorgeschlagene Zielpositionen:\\ \tt - P(1,4): aktuell belegt mit S\\ \tt - O(2,4): aktuell leer (◌)\\ \tt - T(3,4): aktuell Trennlinie (─), ungültig\\ \tt - S(4,4): aktuell leer (◌)\\ \tt - D(5,4): ungültig (nur bis Spalte 5, aber (5,4) ist Trennlinie? Zeile 4 hat Spalte 4: ◌)\\ \tt - A(1,5): aktuell leer (◌)\\ \tt - M(2,5): aktuell Trennlinie (─), ungültig\\ \tt \\ \tt Korrektur: (5,4) existiert nicht, da Zeile 4 nur Spalten 1-5 hat. (5,4) wäre Zeile 4, Spalte 4? Aber Spalte 4 ist in Zeile 4 leer (◌). \\ \tt Tatsächlich:\\ \tt - (1,4): belegt (S)\\ \tt - (2,4): leer (◌)\\ \tt - (3,4): Trennlinie (─), ungültig\\ \tt - (4,4): leer (◌)\\ \tt - (5,4): nicht existent (nur Spalten 1-5)\\ \tt - (1,5): leer (◌)\\ \tt - (2,5): Trennlinie (─), ungültig\\ \tt \\ \tt Daher sind (3,4), (5,4), (2,5) ungültig. Ich muss dies mitteilen.\\ \tt \\ \tt SPRICH: (1,4) ist belegt, (3,4) und (2,5) sind Trennlinien, (5,4) existiert nicht. Bitte alternative Positionen vorschlagen, z.B. (4,4) für P, (1,5) für O, (3,5) für T, (4,5) für S, (5,5) für D, (1,6) für A, (2,6) für M. 
	  } 
	   } 
	   } 
	 & & \\ 
 

    \theutterance \stepcounter{utterance}  

    & & & \multicolumn{2}{p{0.3\linewidth}}{\cellcolor[rgb]{0.95,0.95,0.95}{%
	\makecell[{{p{\linewidth}}}]{% 
	  \tt {\tiny [GM$|$GM]}  
	 Invalid format: head is not empty: 'Prüfe die vorgeschlagenen Positionen auf meinem aktuellen Raster:\\ \tt \\ \tt Aktuelles Raster:\\ \tt \\ \tt  1234567\\ \tt ╔═══╤═╤═╗\\ \tt ║PM◌│S│◌║ 1   // (1,1): P, (1,2): M, (1,3): ◌, (1,4): S, (1,5): ◌\\ \tt ║DAO│◌├─╢ 2   // (2,1): D, (2,2): A, (2,3): O, (2,4): ◌, (2,5): ─ (Trennlinie)\\ \tt ║◌◌T├─┘◌║ 3   // (3,1): ◌, (3,2): ◌, (3,3): T, (3,4): ─, (3,5): ◌\\ \tt ║◌◌◌│◌◌◌║ 4   // (4,1): ◌, (4,2): ◌, (4,3): ◌, (4,4): ◌, (4,5): ◌\\ \tt ╟─┬─┼───╢ 5   // Trennlinie\\ \tt ║◌│◌│◌◌◌║ 6   // (6,1): ◌, (6,2): ◌, (6,3): ◌, (6,4): ◌, (6,5): ◌\\ \tt ║◌│◌│◌◌◌║ 7   // (7,1): ◌, (7,2): ◌, (7,3): ◌, (7,4): ◌, (7,5): ◌\\ \tt ╚═╧═╧═══╝\\ \tt \\ \tt \\ \tt Vorgeschlagene Zielpositionen:\\ \tt - P(1,4): aktuell belegt mit S\\ \tt - O(2,4): aktuell leer (◌)\\ \tt - T(3,4): aktuell Trennlinie (─), ungültig\\ \tt - S(4,4): aktuell leer (◌)\\ \tt - D(5,4): ungültig (nur bis Spalte 5, aber (5,4) ist Trennlinie? Zeile 4 hat Spalte 4: ◌)\\ \tt - A(1,5): aktuell leer (◌)\\ \tt - M(2,5): aktuell Trennlinie (─), ungültig\\ \tt \\ \tt Korrektur: (5,4) existiert nicht, da Zeile 4 nur Spalten 1-5 hat. (5,4) wäre Zeile 4, Spalte 4? Aber Spalte 4 ist in Zeile 4 leer (◌). \\ \tt Tatsächlich:\\ \tt - (1,4): belegt (S)\\ \tt - (2,4): leer (◌)\\ \tt - (3,4): Trennlinie (─), ungültig\\ \tt - (4,4): leer (◌)\\ \tt - (5,4): nicht existent (nur Spalten 1-5)\\ \tt - (1,5): leer (◌)\\ \tt - (2,5): Trennlinie (─), ungültig\\ \tt \\ \tt Daher sind (3,4), (5,4), (2,5) ungültig. Ich muss dies mitteilen.\\ \tt \\ \tt ' 
	  } 
	   } 
	   } 
	 & & \\ 
 

    \theutterance \stepcounter{utterance}  

    & & & \multicolumn{2}{p{0.3\linewidth}}{\cellcolor[rgb]{0.95,0.95,0.95}{%
	\makecell[{{p{\linewidth}}}]{% 
	  \tt {\tiny [GM$|$GM]}  
	 Maximum number of penalties exceeded 
	  } 
	   } 
	   } 
	 & & \\ 
 

    \theutterance \stepcounter{utterance}  

    & & & \multicolumn{2}{p{0.3\linewidth}}{\cellcolor[rgb]{0.95,0.95,0.95}{%
	\makecell[{{p{\linewidth}}}]{% 
	  \tt {\tiny [GM$|$GM]}  
	 * success: False\\ \tt * lose: True\\ \tt * aborted: True\\ \tt -------\\ \tt * Shifts: 8.00\\ \tt * Max Shifts: 12.00\\ \tt * Min Shifts: 6.00\\ \tt * End Distance Sum: 18.20\\ \tt * Init Distance Sum: 30.19\\ \tt * Expected Distance Sum: 29.33\\ \tt * Penalties: 17.00\\ \tt * Max Penalties: 16.00\\ \tt * Rounds: 18.00\\ \tt * Max Rounds: 28.00\\ \tt * Object Count: 7.00\\ \tt  
	  } 
	   } 
	   } 
	 & & \\ 
 

    \theutterance \stepcounter{utterance}  

    & & & \multicolumn{2}{p{0.3\linewidth}}{\cellcolor[rgb]{0.95,0.95,0.95}{%
	\makecell[{{p{\linewidth}}}]{% 
	  \tt {\tiny [GM$|$GM]}  
	 * Main Score: 0.00\\ \tt \\ \tt -------\\ \tt * Distance Score: 0.39\\ \tt * Consistency Score: 0.71\\ \tt * Coverage Score: 0.66\\ \tt * Penalty Score: 0.00\\ \tt * Alternative Penalty Score: 0.00\\ \tt * Alternative Main Score: 0.00\\ \tt \\ \tt -------\\ \tt * Shifts: 8.00\\ \tt * Max Shifts: 12.00\\ \tt * Min Shifts: 6.00\\ \tt * End Distance Sum: 18.20\\ \tt * Init Distance Sum: 30.19\\ \tt * Expected Distance Sum: 29.33\\ \tt * Penalties: 17.00\\ \tt * Max Penalties: 16.00\\ \tt * Rounds: 18.00\\ \tt * Max Rounds: 28.00\\ \tt * Object Count: 7.00\\ \tt  
	  } 
	   } 
	   } 
	 & & \\ 
 

\end{supertabular}
}

\end{document}
